\chapter{Fubini--Study, PL, and psh functions}
\index{Fubini--Study function}\index{piecewise linear function}\index{plurisubharmonic function!BJ}
\label{chap:bj-fs-pl-psh}

Four classes of functions occur in the trivially valued theory and must not
be conflated.  Ordinary PL functions are bounded continuous test functions.
Fubini--Study functions form positive cones indexed by line bundles.
Generically finite Fubini--Study functions may have poles on
proper-support points.  Finally, psh functions are decreasing limits of the
generically finite class, with the numerical classes also allowed to vary.

This chapter gives the definitions and the basic closure theorem of
Boucksom--Jonsson.  No envelope theorem is used.  The distinction matters:
closure under decreasing limits is built into the theory, whereas closure
under arbitrary bounded suprema is precisely the difficult question of
\Cref{chap:bj-envelope-history}.

\section{Orientation}
\label{sec:ch34-orientation}

\subsection{What has to be built}
\label{sec:ch34-what}

\Cref{chap:trivial-analytification} produced a compact Hausdorff space
\(X^{\An}\) of semivaluations, with a scaling action, a dense set
\(X^{\divis}\) of divisorial points, and a supply of continuous functions
\(v\mapsto\log|\mathfrak a|(v)=-v(\mathfrak a)\) attached to coherent
ideals.  \Cref{chap:bj-test-configurations} produced degenerations, flag
ideals, graded norms, and the analytic function
\eqref{eq:norm-fs-function} attached to a filtration.  What is still
missing is the class of \emph{potentials}: the objects on which energies,
Monge--Amp\`ere measures, and the envelope theorem will operate.

Over a nontrivially valued field this class was built in
\Cref{chap:bfj-psh} out of model functions, by imposing a positivity
condition on every model and passing to decreasing limits.  Over a trivially
valued field the same programme runs, but with two structural differences
that make the present chapter genuinely different rather than a translation.

\emph{First, the value group of the ground field is trivial.}  A model
metric over a discretely valued field is built out of sections and elements
of \(K^\times\); \Cref{prop:metrics-fs-model} absorbed the weights of a
Fubini--Study metric into the sections precisely because
\(v(K^\times)\ne0\).  When \(K=k\) is trivially valued this device is
unavailable, and \Cref{prop:metrics-fs-real-weights} fails: real weights are
not limits of weights in the value group, because the value group is
\(\{0\}\).  Fubini--Study functions with prescribed real weights are
therefore \emph{primitive} objects here, not derived ones.  That is the
starting point announced at the end of \Cref{sec:metrics-fs}.

\emph{Second, the analytic space contains points of proper support.}  If
\(v\in X^{\An}\) has support \(Z=s(v)\subsetneq X\) and a section \(s\)
vanishes identically along \(Z\), then \(v(s)=+\infty\) and
\(\log|s|(v)=-\infty\).  Such points are not removable: they are the
scaling limits \(\lim_{t\to+\infty}tv\) of ordinary valuations
\eqref{eq:valuation-scaling}, they carry the whole Zariski topology of
\(X\) inside \(X^{\An}\), and they are the points at which the envelope
theorem of \Cref{chap:bj-envelope-statement} is hardest to state correctly.
Every definition below is therefore formulated for functions with values in
\(\mathbf R\cup\{-\infty\}\), and every operation is checked against the
extended-real ledger of \Cref{sec:extended-arithmetic}.

\subsection{Four classes of functions, kept apart}
\label{sec:ch34-four-classes}

The reader should carry the following table through the chapter.  The four
classes are related by inclusions and by limit operations, but they are
built by different mechanisms and have different closure properties.

\begin{center}
\small
\begin{tabular}{p{.17\textwidth}|p{.22\textwidth}|p{.26\textwidth}|p{.24\textwidth}}
\toprule
class & built from & values & closed under\\
\midrule
\(\PL(X)\)
 & flag ideals, \(\mathbf Q\)-linear combinations
 & real, bounded, continuous
 & sums, differences, \(\max\), \(\min\)\\
\(\cH_\Lambda(L)\)
 & base-point-free weighted section systems
 & real, bounded, continuous
 & sums across classes, same-class \(\max\), constants\\
\(\cH^{\mathrm{gf}}_\Lambda(L)\)
 & arbitrary weighted section systems, generically finite
 & \(\mathbf R\cup\{-\infty\}\); poles on proper support
 & sums across classes, same-class \(\max\), constants\\
\(\PSH_X(\theta)\)
 & decreasing nets of the previous class, with varying classes
 & \(\mathbf R\cup\{-\infty\}\); finite on \(X^{\lin}\)
 & generically finite decreasing nets, sums, \(\max\), convex combinations\\
\bottomrule
\end{tabular}
\end{center}

Three inclusions are elementary and will be proved:
\begin{equation}
\label{eq:ch34-hierarchy}
 \begin{aligned}
  \cH_\Lambda(L)
    &=\cH^{\mathrm{gf}}_\Lambda(L)\cap C^0(X^{\An}),\\
  \cH^{\mathrm{gf}}_{\mathbf R}(L)
    &\subset\PSH_X(c_1(L)),\\
  \mathbf Q_{>0}\cH(L)&\subset\PL^+(X).
 \end{aligned}
\end{equation}
The last inclusion holds for every \(L\), and it is an equality when \(L\)
is ample, by \cite[Corollary~2.26]{BJ22}.  Two further relations are
\emph{not} inclusions and must be watched.  A PL function need not be
\(\theta\)-psh for any \(\theta\), because psh functions are decreasing for
the order \eqref{eq:semivaluation-order} while a general PL function is not.
And an element of \(\PL(X)\cap\PSH_X(\theta)\) need not be a Fubini--Study
function for a single line bundle: relative nefness is weaker than relative
semiampleness.

\subsection{Standing hypotheses and notation}
\label{sec:ch34-hypotheses}

Throughout the chapter \(k\) is an algebraically closed field equipped with
the trivial absolute value.  Following \cite[\S1]{BJ22}, and departing in
this part from the default of the conventions sheet, a \emph{variety} over
\(k\) means a reduced separated scheme of finite type; it need not be
irreducible.  Thus \(X\) is a projective \(k\)-variety with irreducible
components \((X_\alpha)_{\alpha\in A}\), a finite family.  We write
\(N^1(X)=N^1(X)_{\mathbf R}\) for the real N\'eron--Severi space and
\(\Pic(X)_{\mathbf Q}=\Pic(X)\otimes\mathbf Q\).  Statements requiring a
function field, a normalization, or a birational model assume in addition
that \(X\) is integral, and say so.

The analytic vocabulary is that of \Cref{chap:trivial-analytification}.  We
write \(s(v)\) for the support of \(v\in X^{\An}\)
(\Cref{def:triv-support}), \(c(v)\) for its center, and
\begin{equation}
\label{eq:ch34-loci}
 X^{\triv}\subset X^{\An},\qquad
 X^{\val}\subset X^{\An},\qquad
 X^{\divis}\subset X^{\val},\qquad
 X^{\lin}\subset X^{\val}
\end{equation}
for, respectively: the trivial semivaluations \(v_{Y,\triv}\) of the
irreducible subvarieties \(Y\subset X\); the semivaluations of full support,
that is, those whose support is an irreducible component
\eqref{eq:triv-full-support-locus}; the divisorial points, namely
\(v_{\triv,\alpha}:=v_{X_\alpha,\triv}\) together with the positive rational
multiples \(t\ord_E\) of prime-divisor valuations over the components; and
the valuations of linear growth, studied systematically in
\Cref{chap:bj-linear-growth}.  A point of
\(X^{\An}\setminus X^{\val}\) is said to have \emph{proper support}.

For a coherent ideal \(\mathfrak a\subset\cO_X\) we use throughout the
notation of \eqref{eq:ideal-evaluation-triv} in the logarithmic form
\begin{equation}
\label{eq:ch34-log-ideal}
 \log|\mathfrak a|:X^{\An}\longrightarrow[-\infty,0],
 \qquad
 \log|\mathfrak a|(v):=-v(\mathfrak a),
\end{equation}
so that \eqref{eq:ideal-tropical-identities} becomes the max-plus dictionary
\begin{equation}
\label{eq:ch34-log-ideal-semiring}
 \log|\mathfrak a\mathfrak b|=\log|\mathfrak a|+\log|\mathfrak b|,
 \qquad
 \log|\mathfrak a+\mathfrak b|
 =\max\{\log|\mathfrak a|,\log|\mathfrak b|\},
\end{equation}
and \(\log|\mathfrak a|(v)=-\infty\) exactly when
\(s(v)\subset V(\mathfrak a)\).  Note that the first identity in
\eqref{eq:ch34-log-ideal-semiring} is an identity in
\(\mathbf R\cup\{-\infty\}\) with no exception: the forbidden sum
\((+\infty)+(-\infty)\) of ledger entry~(2) cannot occur, because both terms
lie in \([-\infty,0]\).

\subsection{Where the value \texorpdfstring{\(-\infty\)}{-infinity} comes from,
and how it is handled}
\label{sec:ch34-minus-infinity}

It is worth stating at the outset exactly which operations of this chapter
can produce an infinite value, and which ledger entry of
\Cref{sec:extended-arithmetic} legitimizes each one.  The list is short, and
it is complete: no other extended-real operation occurs in the chapter.

\begin{enumerate}
\item \emph{Evaluation of a section.}  \(\log|s|(v)=-v(s)\) equals
  \(-\infty\) exactly when \(s\) vanishes identically along \(s(v)\).  This
  is the only source of infinite values; everything else propagates it.
\item \emph{Finite maxima.}  \(\max\) is computed in the complete lattice
  \(\Rbar\), ledger entry~(1).  No hypothesis is needed, and a term equal to
  \(-\infty\) simply never realizes the maximum unless all terms do.
\item \emph{Sums.}  A sum of two functions with values in
  \(\mathbf R\cup\{-\infty\}\) is always defined, by ledger entry~(2); the
  forbidden mixed sum needs a \(+\infty\), which never occurs here.
\item \emph{Positive scaling.}  \(t\varphi\) for \(t>0\) is defined by
  ledger entry~(3).  The product \(0\cdot(-\infty)\) is \emph{not} defined,
  which is why the limit of \(t\varphi\) as \(t\downarrow0\) is stated below
  as a statement about regularized limits and not as an evaluation.
\item \emph{Decreasing limits.}  A decreasing net in
  \(\mathbf R\cup\{-\infty\}\) converges in \(\Rbar\) to its infimum by
  \Cref{lem:monotone-net-converges}; the limit may be \(-\infty\), and the
  definition of a psh function excludes exactly the case in which this
  happens identically on a component.
\item \emph{Suprema and regularization.}  \(\sup_i\varphi_i\) and
  \((\sup_i\varphi_i)^*\) are defined for every family, using ledger
  entry~(5) for the empty family and \eqref{eq:usc-regularization} for the
  star.
\end{enumerate}

Two operations are \emph{never} performed: no difference of two functions
that may both equal \(-\infty\) is formed, and no quotient or product with a
possibly vanishing factor is formed.  Comparisons are therefore always
written in the relative form \(\varphi\le\psi+C\) rather than
\(\varphi-\psi\le C\), in accordance with
\Cref{rem:forbidden-cancellations}.  When a proof below appears to subtract,
it subtracts a \emph{finite} quantity---typically a value at a trivial
semivaluation, which \Cref{prop:ch34-generic-finiteness} shows to be
finite---and the corresponding ledger entry is named on the spot.

\subsection{Roadmap}
\label{sec:ch34-roadmap}

\Cref{sec:ch34-sections} sets up the evaluation of sections on
semivaluations, proves that the resulting functions are continuous with
values in \([-\infty,0]\), and identifies their polar loci.
\Cref{sec:ch34-fs} defines generically finite Fubini--Study functions,
normalizes their presentations, and proves the full operational calculus,
including the two presentation lemmas that the envelope proof consumes.
\Cref{sec:ch34-pl} treats \(\PL(X)\), its lattice structure, the density
theorem, and the correspondence with test configurations.
\Cref{sec:ch34-dual-complexes} realizes PL functions as piecewise affine
functions on dual complexes of SNC test configurations and computes two
examples completely.  \Cref{sec:ch34-pl-positivity} defines relative
nefness at PL level and checks it on those examples.
\Cref{sec:ch34-psh} contains the definition of \(\PSH_X(\theta)\) and the
theorem that it is the minimal class with the two required properties; its
proof is the diagonal net lemma, which is where the failure of first
countability is felt.  \Cref{sec:ch34-closure} proves the closure calculus,
with an explicit extended-real audit.  \Cref{sec:ch34-proper-support}
studies proper support and the value \(-\infty\), and proves the finiteness
and determination statements that \Cref{chap:bj-envelope-statement}
consumes.  \Cref{sec:ch34-compactness} treats compactness modulo constants
and regularization, and marks precisely the point at which the envelope
property is needed and is not available.  The chapter closes with a
dictionary and a source note.

\section{Sections, ideals, and the value \texorpdfstring{\(-\infty\)}{-infinity}}
\label{sec:ch34-sections}

\subsection{Evaluating a section on a semivaluation}

Let \(L\) be a line bundle on \(X\), let \(m\ge1\), and let
\(s\in H^0(X,mL)\).  For \(v\in X^{\An}\) choose a trivialization
\(\tau\) of \(mL\) on a Zariski open neighbourhood of the center \(c(v)\),
write \(s=f\tau\) with \(f\in\cO_{X,c(v)}\), and set
\begin{equation}
\label{eq:section-log-value}
 v(s):=v(f)\in[0,+\infty],
 \qquad
 \log|s|(v):=-v(s)\in[-\infty,0].
\end{equation}

\begin{lemma}
\label{lem:ch34-frame-independence}
The value \eqref{eq:section-log-value} does not depend on the choice of
\(\tau\), and \(v(s)\ge0\) always.  Moreover \(v(s)=0\) if and only if
\(s\) does not vanish at \(c(v)\).
\end{lemma}

\begin{proof}
Two trivializations of \(mL\) near \(c(v)\) differ by a unit
\(u\in\cO_{X,c(v)}^\times\), so \(f\) changes to \(u^{-1}f\).  Both \(u\)
and \(u^{-1}\) lie in \(\cO_{X,c(v)}\), on which \(v\) is nonnegative by the
definition of a center, so \(v(u)\ge0\) and \(v(u^{-1})\ge0\); since
\(v(u)+v(u^{-1})=v(1)=0\) in \([0,+\infty]\), both vanish.  Hence
\(v(u^{-1}f)=v(f)\).  Nonnegativity of \(v(f)\) is again the defining
property of the center.  Finally, the center is characterized by the fact
that its maximal ideal is \(\{g\in\cO_{X,c(v)}\mid v(g)>0\}\); thus
\(v(f)>0\) if and only if \(f\in\mathfrak m_{c(v)}\), that is, if and only
if \(s(c(v))=0\).
\end{proof}

The function \(\log|s|\) is the additive weight of the section \(s\) for the
\emph{trivial metric} on \(L^{\An}\), the metric whose unit ball is the
integral structure of any algebraic trivialization; compare
\Cref{sec:metrics-continuous} and the sign conventions of
\Cref{app:normalizations}.  We shall not need the metric language, but it
explains why \(\log|s|\le0\): the trivial metric has \(\|\tau\|\equiv1\) on
the locus where \(\tau\) is a frame, and a section can only shrink.

Continuity is not formal, because the center map is not continuous.  The
following two lemmas settle it.

\begin{lemma}
\label{lem:ch34-center-anticontinuous}
For every Zariski closed \(Z\subset X\) the set \(c^{-1}(Z)\) is open in
\(X^{\An}\), and for every Zariski open \(U\subset X\) the set
\(c^{-1}(U)\) is closed in \(X^{\An}\) and contained in \(U^{\An}\).
\end{lemma}

\begin{proof}
Write \(Z=V(\mathfrak b)\) for a coherent ideal \(\mathfrak b\).  By the
characterization of the center recalled in the proof of
\Cref{lem:ch34-frame-independence}, we have \(c(v)\in Z\) if and only if
\(\mathfrak b_{c(v)}\subset\mathfrak m_{c(v)}\), that is, if and only if
\(v(\mathfrak b)>0\).  The function \(v\mapsto v(\mathfrak b)\) is
continuous with values in \([0,+\infty]\) by
\Cref{def:trivial-analytification} and the discussion following
\eqref{eq:ideal-evaluation-triv}; hence
\(c^{-1}(Z)=\{v(\mathfrak b)>0\}\) is open, and its complement
\(c^{-1}(X\setminus Z)\) is closed.

For the last assertion let \(U\) be open and \(c(v)\in U\).  The support
\(s(v)\) is the closure of the point \(\xi_v\) and contains \(c(v)\), so
\(U\cap s(v)\) is a nonempty open subset of the irreducible set \(s(v)\) and
therefore contains its generic point \(\xi_v\).  Thus \(v\in U^{\An}\).
\end{proof}

\begin{lemma}
\label{lem:ch34-log-s-continuous}
For every \(s\in H^0(X,mL)\) the function
\(\log|s|:X^{\An}\to[-\infty,0]\) is continuous, where the target carries
the order topology.  Equivalently, \(|s|=\rme^{\log|s|}\) is a continuous
function \(X^{\An}\to[0,1]\).
\end{lemma}

\begin{proof}
Choose a finite affine open cover \(X=\bigcup_{i=1}^NU_i\) trivializing
\(mL\), say by frames \(\tau_i\), and write \(s=f_i\tau_i\) with
\(f_i\in\cO_X(U_i)\).  By \Cref{lem:ch34-center-anticontinuous} the sets
\(c^{-1}(U_i)\) are closed, contained in \(U_i^{\An}\), and cover
\(X^{\An}\), since every \(v\) has a center and the \(U_i\) cover \(X\).  A
function on a topological space that is continuous on each member of a
finite closed cover is continuous.  It therefore suffices to prove that
\(v\mapsto v(f_i)\) is continuous on \(c^{-1}(U_i)\).  On \(U_i^{\An}\)
this is part of \Cref{def:trivial-analytification}: the topology of
\(X^{\An}\) makes \(U_i^{\An}\) open and \(v\mapsto|f_i(v)|=\rme^{-v(f_i)}\)
continuous there.  Restricting to the subset \(c^{-1}(U_i)\) preserves
continuity, and on that subset \(v(f_i)=v(s)\) by
\Cref{lem:ch34-frame-independence}.
\end{proof}

\begin{remark}
\label{rem:ch34-finite-cover}
The finiteness of the cover is used twice: to know that a closed cover
detects continuity, and to know that \(c^{-1}(U_i)\) is closed rather than
merely constructible.  Both are available because \(X\) is a scheme of
finite type; no properness is needed for this lemma, although properness is
what guarantees the existence of the center in the first place
\eqref{eq:triv-full-support-locus}.
\end{remark}

\subsection{Polar loci, monotonicity, and homogeneity}

\begin{lemma}
\label{lem:ch34-log-s-properties}
Let \(s,s'\in H^0(X,mL)\), \(t\in H^0(X,m'M)\), and \(v,w\in X^{\An}\),
\(\lambda>0\).  Then:
\begin{enumerate}
\item\label{itm:ch34-logs:polar}
  \(\log|s|(v)=-\infty\) if and only if \(s\) vanishes identically along
  \(s(v)\); equivalently, if and only if \(s(v)\subset\{s=0\}\);
\item\label{itm:ch34-logs:product}
  \(\log|st|=\log|s|+\log|t|\) as functions with values in
  \(\mathbf R\cup\{-\infty\}\), the sum being always defined;
\item\label{itm:ch34-logs:sum}
  \(\log|s+s'|\le\max\{\log|s|,\log|s'|\}\);
\item\label{itm:ch34-logs:decreasing}
  \(\log|s|\) is decreasing for the order \eqref{eq:semivaluation-order}:
  \(v\le w\) implies \(\log|s|(v)\ge\log|s|(w)\);
\item\label{itm:ch34-logs:homogeneous}
  \(\log|s|(\lambda v)=\lambda\log|s|(v)\), with
  \(\lambda\cdot(-\infty)=-\infty\) as in ledger entry~(3);
\item\label{itm:ch34-logs:max}
  for every irreducible subvariety \(Y\subset X\),
  \(\sup_{Y^{\An}}\log|s|=\log|s|(v_{Y,\triv})\).
\end{enumerate}
\par
\end{lemma}

\begin{proof}
\emph{(1).}  With notation as in
\eqref{eq:section-log-value}, \(v(f)=+\infty\) means that \(f\) lies in the
support prime \(\mathfrak p_v\) of \(v\), that is, that \(f\) vanishes at
the generic point \(\xi_v\) of \(s(v)\); since \(s(v)\) is the closure of
\(\xi_v\), this is equivalent to \(s\) vanishing identically along
\(s(v)\).

\emph{(2).}  Choosing frames of \(mL\) and
\(m'M\) at \(c(v)\), their product is a frame of \(mL+m'M\) and
multiplicativity of \(v\) gives \(v(st)=v(s)+v(t)\) in \([0,+\infty]\).
Negating and using ledger entry~(3) with \(\lambda=1\) yields the identity
in \([-\infty,0]\); no mixed sum occurs because both terms are
\(\le0\).

\emph{(3).}  This is the ultrametric inequality
\(v(f+f')\ge\min\{v(f),v(f')\}\) after negation.

\emph{(4).}  If \(v\le w\) then
\(v(\mathfrak a)\le w(\mathfrak a)\) for every coherent ideal.  The
principal ideal generated by \(f\) near the common center gives
\(v(s)\le w(s)\), whence the claim.  (Two comparable semivaluations have
the same center by \eqref{eq:trivial-endpoints}, so the same frame may be
used for both.)

\emph{(5).}  Immediate from
\((\lambda v)(f)=\lambda v(f)\), which is
\eqref{eq:valuation-scaling}, together with the convention
\(\lambda\cdot(+\infty)=+\infty\).

\emph{(6).}  For \(v\in Y^{\An}\) we have
\(v_{Y,\triv}\le v\) by \eqref{eq:trivial-endpoints}, so
\cref{itm:ch34-logs:decreasing} gives
\(\log|s|(v)\le\log|s|(v_{Y,\triv})\); and \(v_{Y,\triv}\in Y^{\An}\).
\end{proof}

The polar locus in \cref{itm:ch34-logs:polar} deserves a name.

\begin{definition}
\label{def:ch34-analytic-closed-subset}
For a Zariski closed subset \(Z\subset X\) put
\begin{equation}
\label{eq:ch34-Z-an}
 Z^{\An}:=\{v\in X^{\An}\mid s(v)\subset Z\}.
\end{equation}
This is the image of the analytification of \(Z\) with its reduced
structure, and it is a closed subset of \(X^{\An}\); it contains
\(v_{Y,\triv}\) for every irreducible \(Y\subset Z\), and it is disjoint
from \(X^{\val}\) unless \(Z\) contains an irreducible component of \(X\).
\end{definition}

Thus \(\{\log|s|=-\infty\}=\{s=0\}^{\An}\).  In particular, if
\(s\) is a \emph{regular} section, that is, if \(s\) does not vanish
identically on any irreducible component of \(X\), then the polar locus of
\(\log|s|\) is contained in \(X^{\An}\setminus X^{\val}\) and
\(\log|s|\) is finite on the whole of \(X^{\val}\).

\begin{lemma}
\label{lem:ch34-regular-section}
Let \(L\) be a line bundle on \(X\).  The following are equivalent:
\begin{enumerate}
\item \(H^0(X,L)\) contains a regular section;
\item the base locus \(\operatorname{Bs}(L)\) contains no irreducible
  component of \(X\), that is, it is nowhere dense.
\end{enumerate}
\par
\end{lemma}

\begin{proof}
For each component \(X_\alpha\) put
\(V_\alpha=\{s\in H^0(X,L)\mid s|_{X_\alpha}\equiv0\}\), a linear subspace.
A section is regular exactly when it avoids \(\bigcup_\alpha V_\alpha\).
Since \(k\) is algebraically closed it is infinite, and a vector space over
an infinite field is not the union of finitely many proper subspaces; hence
a regular section exists if and only if every \(V_\alpha\) is a proper
subspace.  That \(V_\alpha\ne H^0(X,L)\) says exactly that some section does
not vanish identically on \(X_\alpha\), which is precisely the statement
that \(\operatorname{Bs}(L)\not\supset X_\alpha\).
\end{proof}

Following the standard terminology, a \(\mathbf Q\)-line bundle \(L\) is
\emph{effective} if its asymptotic base locus
\(\mathbf B(L)=\operatorname{Bs}(mL)\), for \(m\) sufficiently divisible, is
nowhere dense, and \emph{semiample} if \(\mathbf B(L)=\varnothing\).  By
\Cref{lem:ch34-regular-section}, \(L\) is effective if and only if some
\(mL\) has a regular section, and then \(m^{-1}\operatorname{div}(s)\) is an
effective \(\mathbf Q\)-Cartier divisor in the class of \(L\); hence
\begin{equation}
\label{eq:ch34-effective-implies-psef}
 L\ \text{effective}\ \Longrightarrow\ c_1(L)\in\operatorname{Psef}(X),
 \qquad
 L\ \text{semiample}\ \Longrightarrow\ c_1(L)\in\Nef(X).
\end{equation}

\begin{example}
\label{ex:ch34-log-s-p1}
Let \(X=\mathbf P^1_k\) with homogeneous coordinates \([Z_0:Z_1]\), let
\(L=\cO(1)\), and let \(s_p\in H^0(X,L)\) be the section vanishing exactly
at a closed point \(p\).  By \Cref{ex:triv-p1-tree}, \(X^{\An}\) is the cone
over \(\mathbf P^1(k)\) with vertex \(v_{X,\triv}\): every point is
\(t\ord_q\) for a unique closed point \(q\) and \(t\in[0,+\infty]\), with
\(0\cdot\ord_q=v_{X,\triv}\) and \(\infty\cdot\ord_q:=v_{q,\triv}\).  Then
\[
 \log|s_p|(t\ord_q)=
 \begin{cases}
 -t,&q=p,\\
 0,&q\ne p,
 \end{cases}
 \qquad
 \log|s_p|(v_{q,\triv})=
 \begin{cases}
 -\infty,&q=p,\\
 0,&q\ne p .
 \end{cases}
\]
So \(\log|s_p|\) is finite on \(X^{\val}\), equal to \(0\) off one ray,
affine of slope \(-1\) along the ray of \(p\), and has a single pole, at the
proper-support point \(v_{p,\triv}\).  Continuity at that pole is visible:
\(\log|s_p|(t\ord_p)=-t\to-\infty\) as \(t\to+\infty\), and
\(t\ord_p\to v_{p,\triv}\).  The polar locus is
\(\{s_p=0\}^{\An}=\{v_{p,\triv}\}\), a single point.
\end{example}

\begin{example}
\label{ex:ch34-log-s-p2}
Let \(X=\mathbf P^2_k\) with coordinates \([x_0:x_1:x_2]\), \(L=\cO(1)\),
and \(s=x_0\), so that \(\{s=0\}=H\) is a line.  Then
\(\{\log|x_0|=-\infty\}=H^{\An}\) is one-dimensional in the Zariski sense
and contains uncountably many points of \(X^{\An}\), among them
\(v_{H,\triv}\) and \(v_{q,\triv}\) for every closed point \(q\in H\).
Concretely, fix \(q=[0:0:1]\) and the local coordinates
\(u=x_0/x_2\), \(w=x_1/x_2\) at \(q\).  For
\(\alpha=(\alpha_1,\alpha_2)\in\mathbf R_{\ge0}^2\) let \(v_\alpha\) be the
monomial valuation of \eqref{eq:quasimonomial-valuation} with
\(v_\alpha(u)=\alpha_1\), \(v_\alpha(w)=\alpha_2\).  Then
\[
 \log|x_0|(v_\alpha)=-\alpha_1 .
\]
Thus \(\log|x_0|\) is finite, and linear, on the whole two-dimensional
quasimonomial cone at \(q\), and tends to \(-\infty\) exactly along the
directions with \(\alpha_1>0\).  Its restriction to the divisorial ray
\(t\ord_H\), \(t\ge0\), is \(-t\); and
\(\lim_{t\to+\infty}t\ord_H=v_{H,\triv}\), where the value is \(-\infty\).
This is the smallest example in which the polar locus of a Fubini--Study
datum is a positive-dimensional analytic set, and it is the example to keep
in mind in \Cref{sec:ch34-proper-support}.
\end{example}

\section{Fubini--Study presentations}
\label{sec:ch34-fs}

\subsection{The definition}

For a \(\mathbf Q\)-line bundle \(L\), choose \(m\geq1\) such that \(mL\)
is an honest line bundle.  A \emph{weighted section system} for \(L\) is a
triple \((m,(s_j)_{1\le j\le r},(\lambda_j)_{1\le j\le r})\) with
\(m\in\mathbf Z_{>0}\), \(0\ne s_j\in H^0(X,mL)\) and
\(\lambda_j\in\mathbf R\).  Its \emph{Fubini--Study function} is
\begin{equation}
\label{eq:gf-fs-presentation}
 \phi=\frac1m\max_{1\leq j\leq r}\{\log|s_j|+\lambda_j\}
 :X^{\An}\longrightarrow\mathbf R\cup\{-\infty\}.
\end{equation}
The maximum is a finite maximum in the complete lattice \(\Rbar\), so
\eqref{eq:gf-fs-presentation} is defined at every point without any
hypothesis, by ledger entry~(1).

\begin{definition}
\label{def:gf-fs-functions}
Let \(\Lambda\subset\mathbf R\) be an additive subgroup.  A
\(\Lambda\)-rational generically finite Fubini--Study function for \(L\) is
a function \(\phi\) of the form \eqref{eq:gf-fs-presentation} with
\(\lambda_j\in\Lambda\), provided that the \(s_j\) do not all vanish
identically on any irreducible component.  The space of such functions is
denoted \(\cH^{\mathrm{gf}}_\Lambda(L)\).

If the \(s_j\) have no common zero on \(X\), then \(\phi\) is finite-valued
on all of \(X^{\An}\), and is called a Fubini--Study function; these form
\(\cH_\Lambda(L)\).
\end{definition}

We write \(\cH(L):=\cH_{\mathbf Q}(L)\) and
\(\cH^{\mathrm{gf}}(L):=\cH^{\mathrm{gf}}_{\mathbf Q}(L)\).  The reader
should compare \eqref{eq:gf-fs-presentation} with the Fubini--Study weight
of \Cref{sec:metrics-fs}: the formula is the same, the normalization
\(m^{-1}\) and the additive placement of \(\lambda_j\) are those fixed in
\Cref{app:normalizations}, and the only differences are that the base-point
condition has been dropped and that the weights are no longer required to
lie in the value group of the ground field---which here is \(\{0\}\).

\subsection{Generic finiteness}

\begin{proposition}
\label{prop:ch34-generic-finiteness}
Let \(\phi\) be given by \eqref{eq:gf-fs-presentation}, and set
\(Z:=\bigcap_j\{s_j=0\}\).  Then:
\begin{enumerate}
\item\label{itm:ch34-gf:cont}
  \(\phi\) is continuous as a map \(X^{\An}\to[-\infty,+\infty)\), and it is
  bounded above by \(m^{-1}\max_j\lambda_j\);
\item\label{itm:ch34-gf:polar}
  \(\{\phi=-\infty\}=Z^{\An}\);
\item\label{itm:ch34-gf:decreasing}
  \(\phi\) is decreasing for \eqref{eq:semivaluation-order} and satisfies
  the maximum principle
  \(\sup_{Y^{\An}}\phi=\phi(v_{Y,\triv})\) for every irreducible
  \(Y\subset X\); in particular
  \(\sup_{X^{\An}}\phi=\max_\alpha\phi(v_{\triv,\alpha})\);
\item\label{itm:ch34-gf:equiv}
  the following are equivalent:
  \begin{enumerate}
  \item \(\phi\) is generically finite, that is,
    \(\phi|_{X_\alpha^{\An}}\not\equiv-\infty\) for every \(\alpha\);
  \item \(\phi(v_{\triv,\alpha})>-\infty\) for every \(\alpha\);
  \item \(\phi\) is finite valued on \(X^{\val}\);
  \item \(Z\) contains no irreducible component of \(X\).
  \end{enumerate}
\item\label{itm:ch34-gf:convex}
  for every \(v\in X^{\An}\), the function
  \(t\mapsto\phi(tv)\) is convex and decreasing on \([0,+\infty)\).
\end{enumerate}
\par
\end{proposition}

\begin{proof}
\emph{(1).}  Each \(\log|s_j|\) is continuous into
\([-\infty,0]\) by \Cref{lem:ch34-log-s-continuous}; adding the constant
\(\lambda_j\) and taking a finite maximum preserves continuity into
\([-\infty,+\infty)\), and \(\log|s_j|\le0\) gives the upper bound.

\emph{(2).}  A finite maximum is \(-\infty\) exactly
when all terms are, so
\(\{\phi=-\infty\}=\bigcap_j\{s_j=0\}^{\An}=Z^{\An}\) by
\Cref{lem:ch34-log-s-properties}\,\cref{itm:ch34-logs:polar} and
\eqref{eq:ch34-Z-an}.

\emph{(3).}  A maximum of decreasing functions
is decreasing, and \Cref{lem:ch34-log-s-properties}\,\cref{itm:ch34-logs:max}
gives the maximum principle after taking the same finite maximum; the last
statement follows because \(X^{\An}=\bigcup_\alpha X_\alpha^{\An}\).

\emph{(4).}  (a)\(\Rightarrow\)(b) is
\cref{itm:ch34-gf:decreasing}, the supremum over \(X_\alpha^{\An}\) being
attained at \(v_{\triv,\alpha}\).  (b)\(\Rightarrow\)(d): by
\cref{itm:ch34-gf:polar}, \(\phi(v_{\triv,\alpha})>-\infty\) means
\(X_\alpha\not\subset Z\).  (d)\(\Rightarrow\)(c): let \(v\in X^{\val}\),
say \(s(v)=X_\alpha\).  Since \(X_\alpha\not\subset Z\), some \(s_j\) does
not vanish identically on \(X_\alpha\), so
\(\log|s_j|(v)>-\infty\) and hence \(\phi(v)>-\infty\).
(c)\(\Rightarrow\)(a) is trivial, since
\(v_{\triv,\alpha}\in X^{\val}\cap X_\alpha^{\An}\).

\emph{(5).}  By
\Cref{lem:ch34-log-s-properties}\,\cref{itm:ch34-logs:homogeneous},
\[
 \phi(tv)=\frac1m\max_j\{t\log|s_j|(v)+\lambda_j\},
\]
a finite maximum of affine functions of \(t\) with nonpositive slopes, where
a slope \(-\infty\) is read as the function equal to \(\lambda_j/m\) at
\(t=0\) and \(-\infty\) for \(t>0\), which is again convex and decreasing in
the extended sense.  A maximum of convex decreasing functions is convex and
decreasing.
\end{proof}

Assertion \cref{itm:ch34-gf:equiv} is why the two conditions in
\Cref{def:gf-fs-functions} are the right ones: the first is equivalent to
finiteness at the trivial valuation of each component, and to finiteness on
the full-support locus \(X^{\val}\).  The second is equivalent to global
generation of the ideal-weighted section system, hence to finiteness
everywhere.  Thus
\begin{equation}
\label{eq:fs-gf-continuous-intersection}
 \cH_\Lambda(L)=
 \cH^{\mathrm{gf}}_\Lambda(L)\cap C^0(X^{\An}).
\end{equation}
Indeed, \(\subset\) is clear from
\Cref{prop:ch34-generic-finiteness}\,\cref{itm:ch34-gf:cont}, and
conversely a generically finite \(\phi\) that is real valued has
\(Z^{\An}=\varnothing\); since \(v_{q,\triv}\in Z^{\An}\) for every closed
point \(q\in Z\), this forces \(Z=\varnothing\).

\begin{lemma}
\label{lem:ch34-veronese}
Let \(\phi\) be as in \eqref{eq:gf-fs-presentation} and let
\(r\in\mathbf Z_{>0}\).  Then
\begin{equation}
\label{eq:ch34-veronese}
 \phi=\frac1{rm}\max_j\{\log|s_j^{\,r}|+r\lambda_j\}.
\end{equation}
Consequently
\(\cH^{\mathrm{gf}}_\Lambda(L)=\cH^{\mathrm{gf}}_{\mathbf Q\Lambda}(L)\) and
\(\cH_\Lambda(L)=\cH_{\mathbf Q\Lambda}(L)\): the weight group may always be
assumed divisible.  Moreover, any two Fubini--Study functions may be
presented with the same integer \(m\).
\end{lemma}

\begin{proof}
Since \(X\) is reduced and \(s_j\ne0\), the power \(s_j^{\,r}\) is a nonzero
section of \(rmL\), and
\Cref{lem:ch34-log-s-properties}\,\cref{itm:ch34-logs:product} gives
\(\log|s_j^{\,r}|=r\log|s_j|\).  Dividing by \(rm\) yields
\eqref{eq:ch34-veronese}, and the vanishing loci are unchanged, so the
generic-finiteness and base-point conditions are unaffected.  For the
divisibility statement, if \(\lambda_j\in\mathbf Q\Lambda\) choose \(r\)
with \(r\lambda_j\in\Lambda\) for all \(j\); \eqref{eq:ch34-veronese}
exhibits \(\phi\) as a \(\Lambda\)-rational presentation.  Given two
presentations with integers \(m\) and \(m'\), replace \(r\) by \(m'\) in the
first and by \(m\) in the second.
\end{proof}

\begin{lemma}
\label{lem:ch34-generic-constancy}
Assume \(X\) integral, and let \(\phi\in\cH^{\mathrm{gf}}_{\mathbf R}(L)\)
be presented as in \eqref{eq:gf-fs-presentation}.  Then
\(U:=X\setminus\bigcup_j\{s_j=0\}\) is a nonempty Zariski open set, and
\begin{equation}
\label{eq:ch34-generic-constancy}
 \phi\equiv\frac1m\max_j\lambda_j=\sup_{X^{\An}}\phi
 \qquad\text{on }c^{-1}(U).
\end{equation}
In general, the same holds componentwise after discarding from the
presentation those \(s_j\) that vanish identically on the component
considered.
\end{lemma}

\begin{proof}
Each \(\{s_j=0\}\) is a proper closed subset of the integral scheme \(X\),
since \(s_j\ne0\); hence \(U\ne\varnothing\).  If \(c(v)\in U\) then no
\(s_j\) vanishes at \(c(v)\), so \(v(s_j)=0\) for all \(j\) by
\Cref{lem:ch34-frame-independence}, and
\(\phi(v)=m^{-1}\max_j\lambda_j\).  As
\(v_{X,\triv}\in c^{-1}(U)\) --- its center is the generic point of \(X\),
which lies in \(U\) --- the value computed is
\(\phi(v_{X,\triv})=\sup\phi\) by
\Cref{prop:ch34-generic-finiteness}\,\cref{itm:ch34-gf:decreasing}.
For the last sentence, restrict the presentation to \(X_\alpha\); the
sections vanishing identically on \(X_\alpha\) contribute the constant
function \(-\infty\), which never realizes a maximum some of whose terms are
finite, and generic finiteness guarantees that at least one term survives.
\end{proof}

The set \(c^{-1}(U)\) is closed in \(X^{\An}\) by
\Cref{lem:ch34-center-anticontinuous}, but it has nonempty interior: it
contains \(c^{-1}(\{q\})\) for every closed point \(q\in U\), and the
latter is open, again by \Cref{lem:ch34-center-anticontinuous}.  Thus a
generically finite Fubini--Study function is constant on a substantial part
of the analytification.  This innocuous statement is used in
\Cref{chap:bj-envelope-consequences} and in the proof of
\Cref{lem:ch34-generic-constancy-corollary} below.

\begin{corollary}
\label{lem:ch34-generic-constancy-corollary}
For \(X\) integral and \(\phi\in\cH^{\mathrm{gf}}_{\mathbf R}(L)\) with a
presentation as above, one has
\(\sup_{X^{\An}}\phi=m^{-1}\max_j\lambda_j\).  In particular
\(\phi-\sup\phi\) is again in \(\cH^{\mathrm{gf}}_{\mathbf R}(L)\), and
\(\cH^{\mathrm{gf}}_{\mathbf R}(L)\) is stable under the normalization
\(\phi\mapsto\phi-\sup\phi\).
\end{corollary}

\begin{proof}
The first assertion is \eqref{eq:ch34-generic-constancy}.  For the second,
subtracting the finite constant \(\sup\phi\) replaces the weights
\(\lambda_j\) by \(\lambda_j-\max_i\lambda_i\), which is legitimate by
ledger entry~(2) since \(\sup\phi\in\mathbf R\).
\end{proof}
\subsection{The operational calculus}

\begin{lemma}\label{lem:gf-fs-operations}
Let \(\phi\in\cH^{\mathrm{gf}}_\mathbf R(L)\) and
\(\psi\in\cH^{\mathrm{gf}}_\mathbf R(M)\).
\begin{enumerate}
\item \(\phi+\psi\in\cH^{\mathrm{gf}}_\mathbf R(L+M)\);
\item if \(M=L\), then
      \(\max\{\phi,\psi\}\in\cH^{\mathrm{gf}}_\mathbf R(L)\);
\item \(\phi+c\in\cH^{\mathrm{gf}}_\mathbf R(L)\) for \(c\in\mathbf R\);
\item pullback by a morphism is GF whenever it remains generically finite on
      every component of the source.
\end{enumerate}
The analogous assertions hold for finite-valued Fubini--Study functions.
\end{lemma}

\begin{proof}
By \Cref{lem:ch34-veronese} we may present
\[
 \phi=\frac1m\max_{1\le i\le r}\{\log|s_i|+\lambda_i\},
 \qquad
 \psi=\frac1{m'}\max_{1\le j\le r'}\{\log|t_j|+\mu_j\},
\]
with \(s_i\in H^0(X,mL)\setminus\{0\}\) and
\(t_j\in H^0(X,m'M)\setminus\{0\}\), and after replacing \(s_i\) by
\(s_i^{m'}\) and \(t_j\) by \(t_j^{m}\) we may assume \(m=m'\), at the cost
of replacing \(m\) by \(mm'\) and the weights by \(m'\lambda_i\) and
\(m\mu_j\).  We therefore assume \(m=m'\) from now on.

\emph{(1) Sums.}  All the numbers \(\log|s_i|(v)+\lambda_i\) and
\(\log|t_j|(v)+\mu_j\) lie in \(\mathbf R\cup\{-\infty\}\) and are bounded
above by \(\max_i\lambda_i\) and \(\max_j\mu_j\) respectively.  Hence
\Cref{lem:safe-supremum} applies pointwise and gives
\[
 m(\phi+\psi)(v)
 =\max_{i,j}\bigl\{\log|s_i|(v)+\log|t_j|(v)+\lambda_i+\mu_j\bigr\}.
\]
Note that the sum on the left is defined at every point: both summands lie
in \(\mathbf R\cup\{-\infty\}\), so the forbidden mixed sum of ledger
entry~(2) cannot arise.  By
\Cref{lem:ch34-log-s-properties}\,\cref{itm:ch34-logs:product} the inner sum
equals \(\log|s_it_j|\) whenever \(s_it_j\ne0\) in
\(H^0(X,m(L+M))\); and
if \(s_it_j=0\), then multiplicativity of every \(v\) gives
\(v(s_i)+v(t_j)=v(0)=+\infty\), so that the corresponding term is
identically \(-\infty\) and may be deleted from the maximum.  At least one
term survives, since \((\phi+\psi)(v_{\triv,\alpha})\) is a sum of two real
numbers by
\Cref{prop:ch34-generic-finiteness}\,\cref{itm:ch34-gf:equiv}.  Deleting the
vanishing products therefore leaves a weighted section system for
\(L+M\) presenting \(\phi+\psi\), whose generic finiteness is exactly the
finiteness of \((\phi+\psi)(v_{\triv,\alpha})\) for every \(\alpha\).  If
\(\phi\) and \(\psi\) are finite valued, so is \(\phi+\psi\), so the
Fubini--Study case follows.

\emph{(2) Same-class maxima.}  With \(M=L\) and a common \(m\), concatenate
the two systems:
\[
 \max\{\phi,\psi\}=\frac1m
 \max\bigl\{\log|s_1|+\lambda_1,\ldots,
            \log|t_{r'}|+\mu_{r'}\bigr\},
\]
which is a weighted section system for \(L\).  Its intersection of vanishing
loci is \(\bigl(\bigcap_i\{s_i=0\}\bigr)\cap\bigl(\bigcap_j\{t_j=0\}\bigr)\),
which contains no component because neither factor does; hence the maximum
is generically finite.  If both systems are base-point free, so is their
union.

\emph{(3) Constants.}  Adding \(c\) replaces every \(\lambda_i\) by
\(\lambda_i+mc\).  This is legitimate at points where \(\phi=-\infty\) by
ledger entry~(2), \(c\) being finite.

\emph{(4) Pullback.}  Let \(f:Y\to X\) be a morphism of projective
\(k\)-varieties.  Then \(f^*s_i\in H^0(Y,mf^*L)\) and
\(f^*\log|s_i|=\log|f^*s_i|\), since composing a semivaluation of \(Y\) with
\(f^\sharp\) computes the value of \(s_i\) in a frame pulled back from a
neighbourhood of \(f(c(w))=c(f(w))\).  Hence
\(f^*\phi=m^{-1}\max_i\{\log|f^*s_i|+\lambda_i\}\).  Deleting the indices
with \(f^*s_i=0\), which contribute the constant \(-\infty\), leaves a
weighted section system for \(f^*L\); it is a presentation of a generically
finite function exactly when \(f^*\phi\not\equiv-\infty\) on any component
of \(Y\), which is the stated qualification.  For finite-valued \(\phi\) no
qualification is needed: base-point freeness is preserved by pullback, so
\(f^*\cH_\Lambda(L)\subset\cH_\Lambda(f^*L)\) unconditionally.
\end{proof}

\begin{warning}{Same-class maxima}
The maximum assertion most often used below concerns two functions in the
\emph{same} class \(L\), where concatenating presentations is immediate.
There is no natural operation taking \(\max\{\phi_L,\psi_M\}\) to class
\(L+M\) without auxiliary section choices.  Accordingly, only the
same-class form is used in later proofs.
\end{warning}

The remaining elementary operations are collected next.  Item
\cref{itm:ch34-ops:scaling} identifies the scaling action of
\eqref{eq:function-scaling-action} on presentations, and is the reason for
the prefactor in \eqref{eq:scaling-fs-weights}.

\begin{proposition}
\label{prop:ch34-fs-operations}
Let \(\Lambda\subset\mathbf R\) be an additive subgroup and
\(L,L'\in\Pic(X)_{\mathbf Q}\).
\begin{enumerate}
\item\label{itm:ch34-ops:rational}
  \(\cH^{\mathrm{gf}}_\Lambda(aL)=a\,\cH^{\mathrm{gf}}_\Lambda(L)\) for
  \(a\in\mathbf Q_{>0}\), and similarly for \(\cH_\Lambda\);
\item\label{itm:ch34-ops:scaling}
  for \(t>0\) with \(t\Lambda\subset\Lambda\), the scaling action
  \eqref{eq:function-scaling-action} preserves
  \(\cH^{\mathrm{gf}}_\Lambda(L)\) and \(\cH_\Lambda(L)\), and acts on
  presentations by \(\lambda_j\mapsto t\lambda_j\), which is
  \eqref{eq:scaling-fs-weights};
\item\label{itm:ch34-ops:semiample}
  if \(N\in\Pic(X)_{\mathbf Q}\) is semiample, then the constant \(0\)
  belongs to \(\cH_0(N)\), and consequently
  \[
   \cH^{\mathrm{gf}}_\Lambda(L)\subset\cH^{\mathrm{gf}}_\Lambda(L+N),
   \qquad
   \cH_\Lambda(L)\subset\cH_\Lambda(L+N);
  \]
\item\label{itm:ch34-ops:nonempty}
  \(\cH^{\mathrm{gf}}_{\mathbf R}(L)\ne\varnothing\) if and only if
  \(\cH^{\mathrm{gf}}_0(L)\ne\varnothing\), if and only if \(L\) is
  effective; and \(\cH_{\mathbf R}(L)\ne\varnothing\) if and only if
  \(\cH_0(L)\ne\varnothing\), if and only if \(L\) is semiample.
\end{enumerate}
\par
\end{proposition}

\begin{proof}
\emph{\cref{itm:ch34-ops:rational}.}  Write \(a=p/q\) with
\(p,q\in\mathbf Z_{>0}\), and present
\(\phi\in\cH^{\mathrm{gf}}_\Lambda(L)\) with exponent \(m\), so that
\(s_i\in H^0(X,mL)\).  Then \(s_i^{\,p}\in H^0(X,pmL)\), and
\(pmL=qm\cdot(aL)\) is an honest line bundle, so \(s_i^{\,p}\) is a section
of \(qm(aL)\).  Using
\Cref{lem:ch34-log-s-properties}\,\cref{itm:ch34-logs:product} and
\Cref{lem:extended-real-safe}\,(2),
\[
 a\phi=\frac p{qm}\max_i\{\log|s_i|+\lambda_i\}
 =\frac1{qm}\max_i\{\log|s_i^{\,p}|+p\lambda_i\}
 \in\cH^{\mathrm{gf}}_{\Lambda}(aL),
\]
the membership using \(p\lambda_i\in\Lambda\) and, for the vanishing loci,
\(\{s_i^{\,p}=0\}=\{s_i=0\}\) since \(X\) is reduced.  Applying the same
construction to \(a^{-1}\) gives the reverse inclusion.  The base-point
condition is likewise preserved, so the assertion for \(\cH_\Lambda\)
follows.

\emph{\cref{itm:ch34-ops:scaling}.}  By
\Cref{lem:ch34-log-s-properties}\,\cref{itm:ch34-logs:homogeneous} the
function \(\log|s_j|\) is homogeneous, so
\[
 (t\cdot\phi)(v)=t\phi(t^{-1}v)
 =\frac tm\max_j\{t^{-1}\log|s_j|(v)+\lambda_j\}
 =\frac1m\max_j\{\log|s_j|(v)+t\lambda_j\},
\]
where the second equality uses homogeneity and the third uses
\Cref{lem:extended-real-safe}\,(2), multiplication by \(t>0\) commuting with
finite maxima in \(\Rbar\).  The vanishing loci are unchanged, so both
membership conditions persist, and \(t\lambda_j\in\Lambda\) by hypothesis.

\emph{\cref{itm:ch34-ops:semiample}.}  Choose \(m\) with \(mN\) an honest
globally generated line bundle, and let \(s_1,\ldots,s_r\) generate
\(H^0(X,mN)\).  For \(v\in X^{\An}\), the sections \(s_i\) generate \(mN\)
at \(c(v)\), so in a frame at \(c(v)\) the germs \(f_i\) generate the unit
ideal and \(\min_iv(f_i)=0\).  Hence
\(m^{-1}\max_i\log|s_i|\equiv0\), which is therefore an element of
\(\cH_0(N)\).  Adding it to a given \(\phi\in\cH^{\mathrm{gf}}_\Lambda(L)\)
and using \Cref{lem:gf-fs-operations}\,(1) gives
\(\phi=\phi+0\in\cH^{\mathrm{gf}}_\Lambda(L+N)\).

\emph{\cref{itm:ch34-ops:nonempty}.}  If \(L\) is effective, choose \(m\) and
a regular \(s\in H^0(X,mL)\) by \Cref{lem:ch34-regular-section}; then
\(m^{-1}\log|s|\in\cH^{\mathrm{gf}}_0(L)\).  Conversely, a generically
finite presentation has, for each \(\alpha\), an index \(j\) with
\(s_j|_{X_\alpha}\ne0\); by
\Cref{prop:ch34-generic-finiteness}\,\cref{itm:ch34-gf:equiv} the
intersection \(Z\) of the vanishing loci contains no component, so
\(\operatorname{Bs}(mL)\subset Z\) is nowhere dense and \(L\) is effective.
For the semiample statement, if \(mL\) is globally generated then
\(0\in\cH_0(L)\) as in \cref{itm:ch34-ops:semiample}; conversely a
finite-valued presentation has \(Z=\varnothing\), so the \(s_j\) generate
\(mL\) at every point and \(mL\) is globally generated.  Finally,
\(\cH_0\subset\cH_{\mathbf R}\) always, which gives the two remaining
implications.
\end{proof}

\begin{remark}
\label{rem:ch34-why-gf}
It is tempting to work only with the finite-valued class \(\cH_\Lambda(L)\),
which consists of continuous functions and is closed under all the
operations above.  This is impossible for two reasons.  First, by
\Cref{prop:ch34-fs-operations}\,\cref{itm:ch34-ops:nonempty} it is empty
unless \(L\) is semiample, hence unless \(c_1(L)\) is nef; the theory would
then have nothing to say about a merely pseudoeffective class, and the
envelope theorem of \Cref{chap:bj-envelope-statement} is stated for
arbitrary \(\theta\in N^1(X)_{\mathbf R}\).  Second, and more fundamentally,
the class must be closed under decreasing limits, and a decreasing limit of
finite-valued Fubini--Study functions is routinely \(-\infty\) somewhere:
\Cref{ex:ch34-decreasing-to-pole} below exhibits a decreasing
\emph{sequence} in \(\cH(\cO_{\mathbf P^1}(1))\) whose limit is
\(\log|s|\), with polar locus a point.  Retaining the value \(-\infty\) is
therefore not a concession; it is forced by the closure requirement.
\end{remark}

\subsection{Two presentation lemmas}

The following two statements are the technical heart of the
``generically finite presentations'' theme.  Both say that a Fubini--Study
presentation can be transported along a birational modification or divided
by a canonical section, and both are consumed in
\Crefrange{chap:bj-cartier-realization}{chap:bj-reflection}.  In each case
the extended-real safeguard is that the conclusion is stated as an
\emph{additive} identity, never as a difference of two functions that may
both be \(-\infty\).

Recall the setting.  Let \(X\) be integral, let
\(0\ne\mathfrak b\subset\cO_X\) be a coherent ideal, and let
\(\pi:Y\to X\) be the blowup of \(\mathfrak b\), with exceptional Cartier
divisor \(E\) determined by
\(\mathfrak b\cdot\cO_Y=\cO_Y(-E)\).  Denote by
\(s_E\in H^0(Y,\cO_Y(E))\) the canonical section, so that
\(\operatorname{div}(s_E)=E\).

\begin{lemma}
\label{lem:ch34-exceptional-section}
With the above notation,
\begin{equation}
\label{eq:ch34-exceptional-section}
 \pi^*\log|\mathfrak b|=\log|s_E|
 \qquad\text{on }Y^{\An},
\end{equation}
and \(\log|s_E|\in\cH^{\mathrm{gf}}_0(E)\).
\end{lemma}

\begin{proof}
Let \(w\in Y^{\An}\) and let \(v=\pi(w)\in X^{\An}\) be its image, so that
\(v(g)=w(\pi^\sharp g)\) for germs \(g\) at \(c(v)=\pi(c(w))\).  If
\(g_1,\ldots,g_N\) generate \(\mathfrak b\) near \(c(v)\), then their
pullbacks generate \(\mathfrak b\cdot\cO_Y=\cO_Y(-E)\) near \(c(w)\), so
\[
 v(\mathfrak b)=\min_\ell w(\pi^\sharp g_\ell)=w\bigl(\cO_Y(-E)\bigr)=w(f),
\]
where \(f\) is a local equation of \(E\) at \(c(w)\).  In the frame
\(f^{-1}\) of \(\cO_Y(E)\), the canonical section \(s_E\) is written
\(s_E=f\cdot f^{-1}\), so \(w(s_E)=w(f)\).  Combining, \(w(s_E)=v(\mathfrak
b)\), which is \eqref{eq:ch34-exceptional-section} after negation.  Since
\(\mathfrak b\ne0\) and \(Y\) is integral, \(s_E\ne0\), so \(\log|s_E|\)
is a one-term presentation with weight \(0\).
\end{proof}

The next proposition uses the following standard fact about the blowup.  We
record it first because it fixes the exponent that the Veronese
normalization must reach.

\begin{lemma}\label{lem:ch34-blowup-pushforward}
With the above notation, \(\mathfrak b^m\subset\pi_*\cO_Y(-mE)\) for all
\(m\in\mathbf N\), with equality for all \(m\) sufficiently large.
\end{lemma}

\begin{proof}
The blowup is the relative \(\Proj\) of the Rees algebra
\(\mathcal R=\bigoplus_{m\ge0}\mathfrak b^m\), which is generated in degree
one, and \(\cO_Y(-E)=\cO_Y(1)\).  The inclusion
\(\mathfrak b^m=\mathcal R_m\subset\pi_*\cO_Y(m)=\pi_*\cO_Y(-mE)\) is the
canonical map from the graded piece to the sections of the twisting sheaf,
and it is an isomorphism for \(m\gg0\) because \(\mathcal R\) is a finitely
generated \(\cO_X\)-algebra generated in degree one; see
\cite[Exercise~II.5.9]{Har}.
\end{proof}

\begin{proposition}
\label{prop:ch34-blowup-presentation}
Let \(X\) be integral, \(\pi:Y\to X\) the blowup of a nonzero coherent ideal
\(\mathfrak b\), and \(L\in\Pic(X)_{\mathbf Q}\).  Then
\begin{equation}
\label{eq:ch34-blowup-presentation}
 \cH^{\mathrm{gf}}_\Lambda(\pi^*L-E)+\log|s_E|
 \ \subset\ \pi^*\cH^{\mathrm{gf}}_\Lambda(L).
\end{equation}
Explicitly, for every \(\phi\in\cH^{\mathrm{gf}}_\Lambda(\pi^*L-E)\) there
is \(\chi\in\cH^{\mathrm{gf}}_\Lambda(L)\) with
\(\phi+\log|s_E|=\pi^*\chi\) as functions on \(Y^{\An}\).
\end{proposition}

\begin{proof}
Let \(\phi\in\cH^{\mathrm{gf}}_\Lambda(\pi^*L-E)\) and write
\(\phi=m^{-1}\max_i\{\log|s_i|+\lambda_i\}\) with
\(s_i\in H^0\bigl(Y,m(\pi^*L-E)\bigr)\setminus\{0\}\).  By
\Cref{lem:ch34-veronese} we may replace \(m\) by any multiple \(rm\) and
\(s_i\) by \(s_i^{\,r}\); by \Cref{lem:ch34-blowup-pushforward} we may
therefore assume that
\begin{equation}
\label{eq:ch34-blowup-exponent}
 \pi_*\cO_Y(-mE)=\mathfrak b^m .
\end{equation}

Fix \(i\).  The section \(s_is_E^{\,m}\) is a section of \(m\pi^*L\) that
vanishes along \(mE\); equivalently, it is a section of
\(m\pi^*L\otimes\cO_Y(-mE)\).  The projection formula and
\eqref{eq:ch34-blowup-exponent} give
\[
 \begin{aligned}
 H^0\bigl(Y,m\pi^*L\otimes\cO_Y(-mE)\bigr)
 &=H^0\bigl(X,mL\otimes\pi_*\cO_Y(-mE)\bigr)\\
 &=H^0\bigl(X,mL\otimes\mathfrak b^m\bigr)
  \subset H^0(X,mL),
 \end{aligned}
\]
so there is \(\sigma_i\in H^0(X,mL)\) with
\(s_is_E^{\,m}=\pi^*\sigma_i\); and \(\sigma_i\ne0\) because
\(s_is_E^{\,m}\ne0\) and \(\pi\) is birational.  Applying
\Cref{lem:ch34-log-s-properties}\,\cref{itm:ch34-logs:product} and
\Cref{lem:safe-supremum} we obtain, on \(Y^{\An}\),
\[
 \begin{aligned}
 \phi+\log|s_E|
 &=\frac1m\max_i\bigl\{\log|s_i|+m\log|s_E|+\lambda_i\bigr\}\\
 &=\frac1m\max_i\bigl\{\log|\pi^*\sigma_i|+\lambda_i\bigr\}
 =\pi^*\chi,\\
 \chi&:=\frac1m\max_i\bigl\{\log|\sigma_i|+\lambda_i\bigr\}.
 \end{aligned}
\]
Here the first equality adds the quantity \(m\log|s_E|(w)\) to a finite
maximum, pointwise in \(w\).  When that quantity is real this is
\Cref{lem:extended-real-safe}\,(1); when it equals \(-\infty\) both sides
equal \(-\infty\), all the terms of the maximum lying in
\(\mathbf R\cup\{-\infty\}\).  In neither case does a forbidden sum occur.

It remains to check that \(\chi\) is generically finite.  Since \(\pi\) is
birational, \(\pi^*\chi(w_{\triv})=\chi(v_{\triv})\) for the trivial
valuations of the two function fields, which are identified.  Now
\(\pi^*\chi=\phi+\log|s_E|\), and both summands are finite at \(w_{\triv}\):
the first because \(\phi\) is generically finite, the second because
\(s_E\ne0\) is a regular section on the integral variety \(Y\), so
\Cref{prop:ch34-generic-finiteness}\,\cref{itm:ch34-gf:equiv} applies.  Hence
\(\chi(v_{\triv})\in\mathbf R\) and
\(\chi\in\cH^{\mathrm{gf}}_\Lambda(L)\).
\end{proof}

\begin{proposition}
\label{prop:ch34-division}
Assume \(X\) integral and normal.  Let \(L\in\Pic(X)_{\mathbf Q}\), let
\(E\) be an effective \(\mathbf Q\)-Cartier divisor on \(X\) with canonical
section \(s_E\), and let \(\phi\in\cH^{\mathrm{gf}}_\Lambda(L)\).  The
following are equivalent:
\begin{enumerate}
\item there is \(C\in\mathbf R\) with \(\phi\le\log|s_E|+C\) on
  \(X^{\An}\);
\item there is \(\psi\in\cH^{\mathrm{gf}}_\Lambda(L-E)\) with
  \(\phi=\psi+\log|s_E|\).
\end{enumerate}
\par
\end{proposition}

\begin{proof}
That (2) implies (1) is immediate with \(C=\sup\psi\), which is finite by
\Cref{lem:ch34-generic-constancy-corollary}.

Assume (1) and write \(\phi=m^{-1}\max_i\{\log|s_i|+\lambda_i\}\) with
\(s_i\in H^0(X,mL)\setminus\{0\}\).  By \Cref{lem:ch34-veronese} we may
enlarge \(m\) so that \(mE\) is an integral Cartier divisor, and we write
\(s_{mE}\in H^0(X,\cO_X(mE))\) for its canonical section, so that
\(\log|s_E|=m^{-1}\log|s_{mE}|\).  For every index \(i\) and every
\(v\in X^{\An}\),
\[
 \log|s_i|(v)+\lambda_i\ \le\ m\phi(v)\ \le\ \log|s_{mE}|(v)+mC,
\]
that is,
\begin{equation}
\label{eq:ch34-division-inequality}
 v(s_i)\ \ge\ v(s_{mE})-mC+\lambda_i
 \qquad\text{in }[0,+\infty].
\end{equation}
Both sides of \eqref{eq:ch34-division-inequality} are read in
\([0,+\infty]\); the inequality is trivially true where the left side is
\(+\infty\), and no subtraction of infinite quantities occurs, the constant
\(mC-\lambda_i\) being real.

Replace \(v\) by \(tv\) for \(t>0\).  Homogeneity
(\Cref{lem:ch34-log-s-properties}\,\cref{itm:ch34-logs:homogeneous}) turns
\eqref{eq:ch34-division-inequality} into
\(t\,v(s_i)\ge t\,v(s_{mE})-mC+\lambda_i\); dividing by \(t\) and letting
\(t\to+\infty\) gives
\[
 v(s_i)\ \ge\ v(s_{mE})\qquad\text{for every }v\in X^{\An}.
\]
Apply this with \(v=\ord_F\) for each prime divisor \(F\) in the support of
\(E\), which is a point of \(X^{\An}\) because \(X\) is normal: it gives
\(\ord_F(s_i)\ge\ord_F(mE)\).  Consider now the rational section
\(\sigma_i:=s_i/s_{mE}\) of the line bundle \(m(L-E)\).  Its associated
Weil divisor is
\[
 \operatorname{div}(\sigma_i)
 =\operatorname{div}(s_i)-mE ,
\]
and this is effective: at a prime divisor \(F\) in the support of \(E\) the
inequality just proved applies, while at any other prime divisor \(G\) one
has \(\ord_G(s_i)\ge0=\ord_G(mE)\).  On a normal variety a rational section
of a line bundle with effective associated Weil divisor is regular, because
\(\cO_X=\bigcap_{F}\cO_{X,F}\) inside \(k(X)\), the intersection running
over the prime divisors.  Hence
\(\sigma_i\in H^0\bigl(X,m(L-E)\bigr)\) and \(s_i=\sigma_is_{mE}\).

Therefore, using
\Cref{lem:ch34-log-s-properties}\,\cref{itm:ch34-logs:product} again,
\[
 \phi=\frac1m\max_i\{\log|\sigma_i|+\log|s_{mE}|+\lambda_i\}
 =\psi+\log|s_E|,
 \qquad
 \psi:=\frac1m\max_i\{\log|\sigma_i|+\lambda_i\}.
\]
Finally \(\psi\) is generically finite: at the trivial valuation
\(v_{\triv}\) both \(\phi(v_{\triv})\) and \(\log|s_E|(v_{\triv})=0\) are
finite, and the displayed identity forces \(\psi(v_{\triv})\) to be finite
as well.  (Here the deduction is a cancellation of the \emph{finite}
quantity \(\log|s_E|(v_{\triv})=0\), which is permitted; had that value been
\(-\infty\), \Cref{rem:forbidden-cancellations} would forbid the step.)
\end{proof}

\begin{remark}
\label{rem:ch34-additive-form}
Statement (2) of \Cref{prop:ch34-division} is usually written
``\(\phi-\log|s_E|\in\cH^{\mathrm{gf}}_\Lambda(L-E)\)''.  On
\(\{s_E=0\}^{\An}\), where \(\log|s_E|=-\infty\), that difference is of the
shape \((-\infty)-(-\infty)\) and is undefined by ledger entry~(4).  The
additive formulation \(\phi=\psi+\log|s_E|\) says the same thing on the
locus where the difference makes sense and remains meaningful elsewhere,
and it is the form in which the statement is used in
\Cref{chap:bj-cartier-realization}.  The same remark applies to
\eqref{eq:ch34-blowup-presentation}.
\end{remark}
\section{Ordinary PL functions}
\label{sec:ch34-pl}

\subsection{The space \texorpdfstring{\(\PL(X)\)}{PL(X)} and its lattice structure}

For a flag ideal \(\mathfrak a\), \Cref{chap:bj-test-configurations}
defined a bounded continuous function \(\phi_{\mathfrak a}\) by
\eqref{eq:flag-ideal-function}; in the weight decomposition
\(\mathfrak a=\sum_\lambda\mathfrak a_\lambda\varpi^{-\lambda}\) this reads
\begin{equation}
\label{eq:ch34-flag-weight-form}
 \phi_{\mathfrak a}
 =\max_{\lambda}\{\log|\mathfrak a_\lambda|+\lambda\},
\end{equation}
a finite maximum because \(\mathfrak a_\lambda=\cO_X\) for
\(\lambda\ll0\) and \(\mathfrak a_\lambda=0\) for \(\lambda\gg0\).

\begin{definition}
\label{def:bj-pl-function}
The space \(\PL(X)\) is the \(\mathbf Q\)-vector space in
\(C^0(X^{\An})\) generated by the functions of flag ideals.  Its real span
is denoted \(\PL_{\mathbf R}(X)\).
\end{definition}

We also set
\begin{equation}
\label{eq:ch34-pl-plus}
 \PL^+(X):=
 \bigl\{m^{-1}\phi_{\mathfrak a}\ \bigm|\
 m\in\mathbf Z_{>0},\ \mathfrak a\ \text{a flag ideal}\bigr\}
 \subset\PL(X).
\end{equation}

\begin{lemma}
\label{lem:ch34-pl-lattice}
The following hold.
\begin{enumerate}
\item\label{itm:ch34-pl:cone}
  \(\PL^+(X)\) is stable under addition, under finite maxima, under
  multiplication by elements of \(\mathbf Q_{>0}\), and under the scaling
  action \eqref{eq:function-scaling-action} of \(\mathbf Q_{>0}\); it
  contains every constant in \(\mathbf Q\).
\item\label{itm:ch34-pl:difference}
  \(\PL(X)=\PL^+(X)-\PL^+(X)\).
\item\label{itm:ch34-pl:lattice}
  \(\PL(X)\) is a \(\mathbf Q\)-vector space stable under finite maxima and
  minima; it contains the rational constants and consists of bounded
  continuous functions taking rational values on \(X^{\divis}\).
\item\label{itm:ch34-pl:realnotlattice}
  \(\PL_{\mathbf R}(X)\) is \emph{not} stable under maxima in general.
\end{enumerate}
\par
\end{lemma}

\begin{proof}
\emph{\cref{itm:ch34-pl:cone}.}  Stability under sums and maxima is
\Cref{lem:flag-semiring-identities} after clearing denominators: for
\(m^{-1}\phi_{\mathfrak a}\) and \((m')^{-1}\phi_{\mathfrak b}\), pass to
the common denominator \(mm'\) and use
\(\phi_{\mathfrak a^{m'}}=m'\phi_{\mathfrak a}\), which is the first
identity of \eqref{eq:flag-semiring} applied \(m'\) times, together with
\[
 \begin{aligned}
  m^{-1}\phi_{\mathfrak a}+(m')^{-1}\phi_{\mathfrak b}
   &=(mm')^{-1}\phi_{\mathfrak a^{m'}\mathfrak b^{m}},\\
  \max\{m^{-1}\phi_{\mathfrak a},(m')^{-1}\phi_{\mathfrak b}\}
   &=(mm')^{-1}\phi_{\mathfrak a^{m'}+\mathfrak b^{m}}.
 \end{aligned}
\]
For \(c=p/q\in\mathbf Q_{>0}\), \(c\,m^{-1}\phi_{\mathfrak a}
=(qm)^{-1}\phi_{\mathfrak a^{p}}\).  The constant \(c\in\mathbf Q\) is
\(\phi_{\varpi^{-c}}\) after clearing denominators, the flag ideal being the
fractional ideal generated by \(\varpi^{-qc}\) with \(m=q\).  Stability
under the scaling action follows from \eqref{eq:ch34-flag-weight-form}: by
\Cref{prop:ch34-fs-operations}\,\cref{itm:ch34-ops:scaling} applied to each
term, \(t\cdot\phi_{\mathfrak a}
=\max_\lambda\{\log|\mathfrak a_\lambda|+t\lambda\}\), and for
\(t\in\mathbf Q_{>0}\) with \(mt\in\mathbf Z\) this is
\(m^{-1}\phi_{\mathfrak a'}\) for the flag ideal
\(\mathfrak a'=\sum_\lambda\mathfrak a_\lambda^{m}\varpi^{-mt\lambda}\).

\emph{\cref{itm:ch34-pl:difference}.}  The inclusion \(\supset\) is clear.
Conversely a general element of \(\PL(X)\) is a \(\mathbf Q\)-linear
combination \(\sum_kc_k\phi_{\mathfrak a_k}\); splitting the coefficients
according to sign and using \cref{itm:ch34-pl:cone} writes it as a
difference of two elements of \(\PL^+(X)\).

\emph{\cref{itm:ch34-pl:lattice}.}  For \(\phi,\psi\in\PL(X)\) write
\(\phi=\phi_1-\phi_2\) and \(\psi=\psi_1-\psi_2\) with all four terms in
\(\PL^+(X)\).  Then
\[
 \max\{\phi,\psi\}
 =\max\{\phi_1+\psi_2,\ \psi_1+\phi_2\}-(\phi_2+\psi_2)
\]
is again a difference of two elements of \(\PL^+(X)\), and
\(\min\{\phi,\psi\}=-\max\{-\phi,-\psi\}\).  Boundedness and continuity
hold for each \(\phi_{\mathfrak a}\) by
\Cref{chap:bj-test-configurations}, and rationality on \(X^{\divis}\)
follows from \eqref{eq:ch34-flag-weight-form}, the value
\(\log|\mathfrak a_\lambda|(t\ord_E)=-t\ord_E(\mathfrak a_\lambda)\) being
rational for \(t\in\mathbf Q_{>0}\).

\emph{\cref{itm:ch34-pl:realnotlattice}.}  This is
\cite[Remark~2.4]{BJ22}.  The obstruction is that taking a maximum after
forming arbitrary real linear combinations can create an irrational wall
that is not induced by any rational model subdivision.  No purported
one-dimensional example is used here: the tempting expression
\(\max\{0,\sqrt2\log|\mathfrak b|\}\) does not exhibit such a wall and is
not a counterexample.
\end{proof}

\begin{example}
\label{ex:ch34-pl-subscheme}
Let \(Z\subset X\) be a closed subscheme with ideal
\(\mathfrak b\subset\cO_X\), and put
\begin{equation}
\label{eq:ch34-phi-Z}
 \phi_Z:=\max\{\log|\mathfrak b|,-1\}\in\PL^+(X),
\end{equation}
which is \(\phi_{\mathfrak a}\) for the flag ideal
\(\mathfrak a=\mathfrak b+(\varpi)\).  Its scalings are
\(t\cdot\phi_Z=\max\{\log|\mathfrak b|,-t\}\) for \(t\in\mathbf Q_{>0}\).
Thus the truncation level is exactly the scaling parameter, and letting
\(t\to+\infty\) recovers the unbounded homogeneous function
\(\log|\mathfrak b|\) of \eqref{eq:positive-homogeneous-pl} as an increasing
limit.  On \(X=\mathbf P^1\) and \(Z=\{p\}\) this is the function equal to
\(\max\{-s,-t\}\) on the ray \(s\ord_p\) and to \(0\) elsewhere, in the
coordinates of \Cref{ex:ch34-log-s-p1}.
\end{example}

\subsection{Density}

\begin{theorem}
\label{thm:ch34-pl-density}
The space \(\PL(X)\) is uniformly dense in \(C^0(X^{\An})\):
\begin{equation}
\label{eq:pl-stone-weierstrass}
 \overline{\PL_{\mathbf R}(X)}^{\,\|\cdot\|_\infty}
 =C^0(X^{\An}).
\end{equation}
\end{theorem}

\begin{proof}
By \Cref{thm:triv-berkovich-compactness} the space \(X^{\An}\) is compact
Hausdorff, so the lattice form of the Stone--Weierstrass theorem applies: a
sublattice of \(C^0(X^{\An})\) containing the constants and separating
points is uniformly dense.  By
\Cref{lem:ch34-pl-lattice}\,\cref{itm:ch34-pl:lattice}, \(\PL(X)\) is a
sublattice containing the rational constants, whose uniform closure
contains all real constants.

It remains to separate points.  Let \(v\ne w\) in \(X^{\An}\).  The topology
of \(X^{\An}\) is generated by the ideal evaluations
\eqref{eq:ideal-evaluation-triv}, and the space is Hausdorff; hence there is
a coherent ideal \(\mathfrak b\) with
\(v(\mathfrak b)\ne w(\mathfrak b)\), that is,
\(\log|\mathfrak b|(v)\ne\log|\mathfrak b|(w)\).  Choose
\(m\in\mathbf Z_{>0}\) larger than \(\min\{-\log|\mathfrak b|(v),
-\log|\mathfrak b|(w)\}\); then the truncation
\(\max\{\log|\mathfrak b|,-m\}\), which lies in \(\PL^+(X)\) by
\Cref{ex:ch34-pl-subscheme} and
\Cref{lem:ch34-pl-lattice}\,\cref{itm:ch34-pl:cone}, still separates \(v\)
from \(w\).  This is \cite[Theorem~2.2 and Lemma~2.3]{BJ22}; the same
argument occurs in the proof of \Cref{thm:triv-divisorial-density}, where it
is combined with the effectivity criterion to prove density of
\(X^{\divis}\).
\end{proof}

\begin{remark}
\label{rem:ch34-truncation}
The unbounded function \(\log|\mathfrak b|\) itself is \emph{not} in
\(\PL(X)\): it takes the value \(-\infty\) on \(V(\mathfrak b)^{\An}\),
whereas \(\PL(X)\) consists of bounded continuous real-valued functions.
The truncation \(\max\{\log|\mathfrak b|,-m\}\) is what enters
\Cref{thm:ch34-pl-density}, and the two are related by the scaling action
of \Cref{ex:ch34-pl-subscheme}.  The homogeneous unbounded objects are
organized separately in \Cref{chap:bj-homogeneity-bdivisors}; see
\eqref{eq:positive-homogeneous-pl}.
\end{remark}

\subsection{Test configurations and Fubini--Study functions}

\begin{proposition}
\label{prop:ch34-fs-flag-ideal}
Let \(L\in\Pic(X)_{\mathbf Q}\) and let
\(\phi:X^{\An}\to\mathbf R\cup\{-\infty\}\).  Then \(\phi\in\cH(L)\) if and
only if there are \(m\in\mathbf Z_{>0}\) with \(mL\) an honest line bundle
and a flag ideal \(\mathfrak a\) on \(X\times\mathbf A^1\) such that
\begin{equation}
\label{eq:ch34-fs-flag-ideal}
 \phi=m^{-1}\phi_{\mathfrak a}
 \qquad\text{and}\qquad
 mL_{\mathrm{triv}}\otimes\mathfrak a\ \text{is globally generated}.
\end{equation}
The second condition is equivalent to global generation of
\(mL\otimes\mathfrak a_\lambda\) for every \(\lambda\).
\end{proposition}

\begin{proof}
Suppose first that \eqref{eq:ch34-fs-flag-ideal} holds.  For each
\(\lambda\), choose sections \(s_{\lambda,1},\ldots,s_{\lambda,r_\lambda}\)
of \(H^0(X,mL\otimes\mathfrak a_\lambda)\subset H^0(X,mL)\) generating
\(mL\otimes\mathfrak a_\lambda\).  For \(v\in X^{\An}\), computing in a
frame of \(mL\) at \(c(v)\), these germs generate
\((\mathfrak a_\lambda)_{c(v)}\), so
\(\min_iv(s_{\lambda,i})=v(\mathfrak a_\lambda)\) and hence
\(\max_i\log|s_{\lambda,i}|=\log|\mathfrak a_\lambda|\).  By
\eqref{eq:ch34-flag-weight-form},
\[
 \phi=\frac1m\max_{\lambda,i}\{\log|s_{\lambda,i}|+\lambda\},
\]
a weighted section system for \(L\) with integer weights.  It is base-point
free: for \(\lambda\) sufficiently negative, \(\mathfrak a_\lambda=\cO_X\), so the
corresponding sections generate \(mL\) and have no common zero.  Thus
\(\phi\in\cH_{\mathbf Z}(L)\subset\cH(L)\).

Conversely, let \(\phi\in\cH(L)\).  After \Cref{lem:ch34-veronese} we may
write
\[
 \phi=m^{-1}\max_i\{\log|s_i|+\lambda_i\},
 \qquad \lambda_i\in\mathbf Z,
 \qquad s_i\in H^0(X,mL),
\]
where the \(s_i\) have no common zero.  For
\(\lambda\in\mathbf Z\) let \(\mathfrak b_\lambda\subset\cO_X\) be the image
of the evaluation map
\[
 \bigoplus_{i:\lambda_i=\lambda}\cO_X(-mL)\longrightarrow\cO_X.
\]
Thus
\(\log|\mathfrak b_\lambda|=\max_{\lambda_i=\lambda}\log|s_i|\), and put
\(\mathfrak a_\lambda:=\sum_{\lambda'\ge\lambda}\mathfrak b_{\lambda'}\).
The \(\mathfrak a_\lambda\) form a decreasing family of ideals which is
\(\cO_X\) for \(\lambda\) small --- because the \(s_i\) have no common zero,
so that \(\sum_\lambda\mathfrak b_\lambda=\cO_X\) --- and zero for
\(\lambda\) large.  Hence
\(\mathfrak a:=\sum_\lambda\mathfrak a_\lambda\varpi^{-\lambda}\) is a flag
ideal, and by \eqref{eq:ch34-flag-weight-form},
\[
 \phi_{\mathfrak a}
 =\max_\lambda\{\log|\mathfrak a_\lambda|+\lambda\}
 =\max_\lambda\max_{\lambda'\ge\lambda}
     \{\log|\mathfrak b_{\lambda'}|+\lambda\}
 =\max_{\lambda'}\{\log|\mathfrak b_{\lambda'}|+\lambda'\}
 =m\phi,
\]
the third equality because \(\lambda\le\lambda'\) makes the term with
\(\lambda=\lambda'\) the largest.  Finally
\(mL\otimes\mathfrak a_\lambda\) is generated by the \(s_i\) with
\(\lambda_i\ge\lambda\), by construction.  This is
\cite[Proposition~2.25]{BJ22}.
\end{proof}

\begin{theorem}\label{thm:tc-pl-correspondence}
For every \(\mathbf Q\)-line bundle \(L\), equivalence classes of test
configurations for \(L\) correspond bijectively to \(\PL(X)\).  Under this
correspondence, semiample test configurations correspond to
\(\cH(L):=\cH_{\mathbf Q}(L)\).  If \(L\) is ample, ample integrally closed
test configurations correspond to \(\cH(L)\).
\end{theorem}

\begin{proof}
\emph{Step 1: the map.}  Let \((\cX,\cL)\) be a test configuration for
\((X,L)\) in the sense of \Cref{def:bj-test-configuration}.  Choose a test
configuration \(\cX'\) with morphisms \(\mu:\cX'\to\cX\) and
\(\rho:\cX'\to X\times\mathbf A^1\); one exists by the directedness
construction of \Cref{prop:ch33-directed-poset}, using the closure of the
graph and flattening, with no resolution hypothesis.  Then
\[
 D:=\mu^*\cL-\rho^*L_{\mathrm{triv}}
\]
is a vertical \(\mathbf Q\)-Cartier divisor class on \(\cX'\), because the
two polarizations agree over \(\mathbf G_m\).  Set
\(\phi_{\cL}:=\phi_D\), the function \(v\mapsto\sigma(v)(D)\) of
\Cref{chap:bj-test-configurations}, where \(\sigma\) is the Gauss section
\eqref{eq:triv-gauss-section}.  Pullback invariance
\(\phi_{\mu^*D}=\phi_D\) shows that \(\phi_{\cL}\) depends only on the
equivalence class of \((\cX,\cL)\), and that the assignment is additive.

\emph{Step 2: bijectivity onto \(\PL(X)\).}  The resulting map from
\(\varinjlim_{\cX}\operatorname{VCar}(\cX)_{\mathbf Q}\) to
\(C^0(X^{\An})\) is injective with image \(\PL(X)\).  This is
\cite[Theorem~2.7]{BJ22}, and we recall its two ingredients, both already
used in \Cref{chap:trivial-analytification}.  Injectivity is the valuative
effectivity criterion \cite[Lemma~2.8]{BJ22}: on an integrally closed
\(\cX\), a vertical \(\mathbf Q\)-Cartier divisor \(D\) is effective if and
only if \(\phi_D(v_E)\ge0\) for every irreducible component \(E\) of
\(\cX_0\); applying this to \(D\) and \(-D\) shows that \(\phi_D=0\) forces
\(D=0\).  Surjectivity onto \(\PL^+(X)\) is \cite[Lemma~2.9]{BJ22}: an
element of \(\PL^+(X)\) is \(m^{-1}\phi_{\mathfrak a}\), and the blowup
\(\mu:\cX\to X\times\mathbf A^1\) of \(\mathfrak a\) has
\(\mathfrak a\cdot\cO_{\cX}=\cO_{\cX}(E)\) with \(E\) vertical Cartier, so
\(m^{-1}\phi_{\mathfrak a}=\phi_{m^{-1}E}\); conversely a
\(\mu\)-semiample vertical divisor is recovered from the flag ideal
\(\mu_*\cO_{\cX}(mD)\).  Together with
\Cref{lem:ch34-pl-lattice}\,\cref{itm:ch34-pl:difference} this gives the
image \(\PL(X)\).

\emph{Step 3: semiample test configurations give Fubini--Study functions.}
Let \((\cX,\cL)\) be semiample.  After replacing \(\cL\) by a multiple and
passing to a higher test configuration we may assume that \(\cL\) is a
globally generated honest line bundle and that \(\cX\) dominates the trivial
test configuration via \(\mu\).  Put \(D=\cL-\mu^*L_{\mathrm{triv}}\), so
that \(\phi_{\cL}=\phi_D\).  Global generation of \(\cL\) makes
\(\cO_{\cX}(D)\) relatively globally generated, so
\(\mathfrak a:=\mu_*\cO_{\cX}(D)\) is a flag ideal with
\(\mathfrak a\cdot\cO_{\cX}=\cO_{\cX}(D)\) and hence
\(\phi_{\mathfrak a}=\phi_D\).  Let \(\mathfrak a_0\subset\mathfrak a\) be
the flag ideal generated by the global sections
\(H^0(X\times\mathbf A^1,L_{\mathrm{triv}}\otimes\mathfrak a)
=H^0(\cX,\cL)\); global generation of \(\cL\) gives
\(\mathfrak a_0\cdot\cO_{\cX}=\cO_{\cX}(D)=\mathfrak a\cdot\cO_{\cX}\),
whence \(\phi_{\mathfrak a_0}=\phi_{\cL}\).  By construction
\(L_{\mathrm{triv}}\otimes\mathfrak a_0\) is globally generated, and
\Cref{prop:ch34-fs-flag-ideal} then gives
\(\phi_{\cL}\in\cH(L)\).

\emph{Step 4: the converse and the ample statement.}  Conversely, let
\(\phi\in\cH(L)\).  By \Cref{prop:ch34-fs-flag-ideal} we may write
\(m\phi=\phi_{\mathfrak a}\) with \(mL_{\mathrm{triv}}\otimes\mathfrak a\)
globally generated.  Let \(\mu:\cX\to X\times\mathbf A^1\) be the blowup of
\(\mathfrak a\), so that \(\mathfrak a\cdot\cO_{\cX}=\cO_{\cX}(D)\) for a
vertical Cartier divisor \(D\); then
\(\cL:=\mu^*(mL)_{\mathrm{triv}}+D\) is globally generated, hence semiample,
and \(\phi_{\cL}=\phi_D=\phi_{\mathfrak a}=m\phi\).  Rescaling by \(m\)
gives a semiample test configuration for \((X,L)\) with associated function
\(\phi\).  Finally, if \(L\) is ample, then every semiample test
configuration is equivalent to a unique ample integrally closed one by
\cite[Proposition~2.30]{BJ22}, which upgrades the bijection of the previous
paragraph to the last assertion.  The two statements just proved are
\cite[Theorem~2.31 and Corollary~2.32]{BJ22}.
\end{proof}

\begin{corollary}
\label{cor:ch34-pl-plus-union}
For every \(L\in\Pic(X)_{\mathbf Q}\) one has
\(\mathbf Q_{>0}\cH(L)\subset\PL^+(X)\), with equality when \(L\) is ample.
Consequently
\begin{equation}
\label{eq:ch34-pl-plus-union}
 \PL^+(X)=\bigcup_{L\ \text{ample}}\cH(L),
\end{equation}
and \(\cH(L)\) spans the \(\mathbf Q\)-vector space \(\PL(X)\) whenever
\(L\) is ample.
\end{corollary}

\begin{proof}
The inclusion \(\cH(L)\subset\PL^+(X)\) is
\Cref{prop:ch34-fs-flag-ideal}, and \(\PL^+(X)\) is stable under
multiplication by \(\mathbf Q_{>0}\) by
\Cref{lem:ch34-pl-lattice}\,\cref{itm:ch34-pl:cone}.

Suppose now that \(L\) is ample and let
\(\phi=m^{-1}\phi_{\mathfrak a}\in\PL^+(X)\).  Only finitely many weight
ideals \(\mathfrak a_\lambda\) are distinct, so Serre's theorem provides
\(r\in\mathbf Z_{>0}\) such that \(rmL\) is an honest line bundle and
\(rmL\otimes\mathfrak a_\lambda\) is globally generated for every
\(\lambda\).  Then
\[
 \frac1r\phi=\frac1{rm}\phi_{\mathfrak a}\in\cH(L)
\]
by \Cref{prop:ch34-fs-flag-ideal}, whence
\(\phi\in\mathbf Q_{>0}\cH(L)\) and
\(\mathbf Q_{>0}\cH(L)=\PL^+(X)\).  Moreover
\Cref{prop:ch34-fs-operations}\,\cref{itm:ch34-ops:rational} gives
\(\phi=r\cdot(r^{-1}\phi)\in\cH(rL)\) with \(rL\) ample, which proves
\eqref{eq:ch34-pl-plus-union}; the reverse inclusion there is the first
sentence.  The last assertion follows from
\Cref{lem:ch34-pl-lattice}\,\cref{itm:ch34-pl:difference}.
\end{proof}

\begin{remark}
\label{rem:ch34-tc-scope}
Three cautions.  First, the correspondence of
\Cref{thm:tc-pl-correspondence} is between \emph{equivalence classes} of
test configurations and PL functions; forgetting either the integral closure
or the linearization destroys injectivity, as recorded in
\Cref{prop:tc-valuative-invariance}.  Second, \(\cH(L)\) corresponds to
\emph{semiample} test configurations, not to nef ones: a relatively nef
polarization need not be relatively semiample, so the class
\(\PL(X)\cap\PSH_X(c_1(L))\) of \Cref{sec:ch34-pl-positivity} is in general
strictly larger than \(\cH(L)\).  Third, nothing here requires resolution of
singularities or characteristic zero: flag ideals are defined for an
arbitrary projective variety over any field.
\end{remark}
\section{Dual complexes and the piecewise affine picture}
\label{sec:ch34-dual-complexes}

The name \emph{piecewise linear} is justified by a geometric picture that we
now make precise.  Nothing in \Cref{def:bj-pl-function} refers to it, and
nothing later depends on its availability; but on the locus where it is
available it converts every statement about vertical divisors into a
statement about affine functions on a compact simplicial complex, and it is
the shortest route to explicit computations.

\subsection{SNC test configurations and their dual complexes}

Throughout this section \(X\) is integral, \(K:=k((\varpi))\), and
\(R:=k[[\varpi]]\), normalized by \(v_K(\varpi)=1\).

\begin{definition}
\label{def:ch34-snc-tc}
A test configuration \(\pi:\cX\to\mathbf A^1\) dominating
\(X\times\mathbf A^1\) is \emph{SNC} if \(\cX\) is regular and the central
fibre
\[
 \cX_0=\sum_{i\in I}b_iE_i
\]
has simple normal crossings support: the \(E_i\) are regular prime divisors
meeting transversally.  We write \(\cX_R:=\cX\times_{\mathbf A^1}\Spec R\)
for the associated \(R\)-model of \(X_K\), obtained by base change along
\(\Spec R\to\mathbf A^1\).
\end{definition}

Base change to \(R\) does not change the special fibre, so \(\cX_R\) is an
SNC model of \(X_K\) in the sense of \Cref{sec:skeleton-snc}, and
\Cref{chap:skeleta} applies verbatim to it.  If \(\operatorname{char}k=0\),
SNC test configurations dominating \(X\times\mathbf A^1\) are cofinal among
all test configurations, by equivariant resolution of singularities
(\Cref{app:birational-geometry}) applied to a \(\mathbf G_m\)-equivariant
compactified model.  In positive characteristic no such cofinality is
available, and this is one of the reasons the definitions of
\Cref{sec:ch34-pl,sec:ch34-psh} are phrased without dual complexes.

\begin{definition}
\label{def:ch34-dual-complex}
Let \(\cX\) be an SNC test configuration.  For \(J\subset I\) with
\(E_J:=\bigcap_{j\in J}E_j\ne\varnothing\), and for
\(\alpha\in\mathbf R_{\ge0}^J\), let \(w_\alpha\) be the monomial valuation
\eqref{eq:quasimonomial-valuation} attached to the stratum \(E_J\) and the
weights \(\alpha\).  Set
\[
 \sigma_J:=\Bigl\{\alpha\in\mathbf R_{\ge0}^J\ \Bigm|\
 \sum_{j\in J}b_j\alpha_j=1\Bigr\},
 \qquad
 \Delta_{\cX_R}:=\bigcup_J\{w_\alpha\mid\alpha\in\sigma_J\}
 \subset X_K^{\An},
\]
the skeleton of \Cref{thm:skeleton-embedding}.  Then
\begin{equation}
\label{eq:ch34-dual-complex}
 \Delta_{\cX}:=\sigma^{-1}\bigl(\Delta_{\cX_R}\bigr)\subset X^{\An},
\end{equation}
where \(\sigma\) is the Gauss section \eqref{eq:triv-gauss-section}.
\end{definition}

\begin{lemma}
\label{lem:ch34-skeleton-in-gauss}
Every \(w_\alpha\) with \(\alpha\in\sigma_J\) is
\(\mathbf G_m(k)\)-invariant and satisfies \(w_\alpha(\varpi)=1\).  Hence
\(\Delta_{\cX_R}\subset\sigma(X^{\An})\), the map \(\sigma\) restricts to a
homeomorphism \(\Delta_{\cX}\xrightarrow{\ \sim\ }\Delta_{\cX_R}\), and
\(\Delta_{\cX}\) is a compact subset of \(X^{\An}\) homeomorphic to the dual
complex of \(\cX_0\).  Its vertices are the divisorial points
\begin{equation}
\label{eq:ch34-dual-vertices}
 v_{E_i}:=b_i^{-1}\ord_{E_i}\big|_{k(X)}\in X^{\divis},
 \qquad i\in I.
\end{equation}
\end{lemma}

\begin{proof}
Locally at the generic point \(\xi_J\) of a connected component of \(E_J\),
one has \(\varpi=u\prod_{j\in J}z_j^{b_j}\) with \(u\) a unit and \(z_j\) a
local equation of \(E_j\), by \Cref{lem:snc-normal-form}.  Since
\(w_\alpha(u)=0\) and \(w_\alpha(z_j)=\alpha_j\), we get
\(w_\alpha(\varpi)=\sum_jb_j\alpha_j=1\).

For invariance, note that each \(E_i\) is a component of \(\cX_0\), hence
\(\mathbf G_m\)-invariant; consequently the stratum \(E_J\) is
\(\mathbf G_m\)-invariant and \(\xi_J\) is a \(\mathbf G_m\)-fixed point of
\(\cX\).  By \Cref{lem:monomial-uniqueness} the valuation \(w_\alpha\)
depends only on the pair \((\xi_J,\alpha)\) and not on the chosen local
equations; therefore \(t\cdot w_\alpha=w_\alpha\) for every
\(t\in\mathbf G_m(k)\).  By the description of the image of \(\sigma\) in
\Cref{chap:trivial-analytification}, \(w_\alpha\in\sigma(X^{\An})\).

Finally \(\sigma\) is a continuous injection between compact Hausdorff
spaces, hence a homeomorphism onto its image, and
\(\Delta_{\cX_R}\) is compact by \Cref{thm:skeleton-embedding}.  The vertex
of \(\sigma_J\) in the direction \(j\) is \(\alpha=b_j^{-1}e_j\), and the
corresponding monomial valuation is \(b_j^{-1}\ord_{E_j}\); restricting to
\(k(X)\subset k(X)(\varpi)\) through \(\sigma\) yields
\eqref{eq:ch34-dual-vertices}, an element of \(X^{\divis}\) --- the trivial
valuation \(v_{X,\triv}\) when \(E_j\) is the strict transform of
\(X\times\{0\}\), and a positive rational multiple of a prime-divisor
valuation otherwise.
\end{proof}

\begin{proposition}
\label{prop:ch34-affine-on-faces}
Let \(\cX\) be an SNC test configuration and let
\(D=\sum_{i\in I}a_iE_i\) be a vertical \(\mathbf R\)-divisor on \(\cX\),
necessarily \(\mathbf R\)-Cartier since \(\cX\) is regular.  Then for every
\(J\) and every \(\alpha\in\sigma_J\),
\begin{equation}
\label{eq:ch34-affine-formula}
 \phi_D\bigl(\sigma^{-1}(w_\alpha)\bigr)=\sum_{j\in J}a_j\alpha_j .
\end{equation}
In particular \(\phi_D\) is affine on every face of \(\Delta_{\cX}\), takes
the value \(a_i/b_i\) at the vertex \(v_{E_i}\), and is determined on
\(\Delta_{\cX}\) by its vertex values.
\end{proposition}

\begin{proof}
By definition \(\phi_D(v)=\sigma(v)(D)\), and for
\(v=\sigma^{-1}(w_\alpha)\) this is \(w_\alpha(D)\).  After multiplying
\(D\) by a positive integer we may assume the \(a_i\) integral, the formula
being homogeneous of degree one in \(D\).  Near \(\xi_J\) the divisor \(D\)
is cut out by \(\prod_{j\in J}z_j^{a_j}\) up to a unit, because the
components \(E_i\) with \(i\notin J\) do not pass through \(\xi_J\).  Hence
\(w_\alpha(D)=\sum_{j\in J}a_j\alpha_j\).  Evaluating at the vertex
\(\alpha=b_i^{-1}e_i\) gives \(a_i/b_i\), and an affine function on a
simplex is determined by its vertex values.  This is
\Cref{lem:model-function-affine}, transported through
\Cref{lem:ch34-skeleton-in-gauss}.
\end{proof}

\begin{proposition}
\label{prop:ch34-retraction}
Let \(\cX\) be an SNC test configuration.  There is a continuous
retraction
\begin{equation}
\label{eq:ch34-retraction}
 r_{\cX}:X^{\An}\longrightarrow\Delta_{\cX},
 \qquad
 r_{\cX}|_{\Delta_{\cX}}=\mathrm{id},\qquad r_{\cX}^2=r_{\cX},
\end{equation}
characterized as follows: for \(v\in X^{\An}\) let \(\xi\) be the center of
\(\sigma(v)\) on \(\cX\), let \(J(\xi)=\{j\mid\xi\in E_j\}\), and let
\(\alpha_j=\sigma(v)(z_j)\) for local equations \(z_j\) of \(E_j\) at
\(\xi\); then \(r_{\cX}(v)=\sigma^{-1}(w_\alpha)\).  Moreover:
\begin{enumerate}
\item\label{itm:ch34-retr:divisor}
  \(\phi_D=\phi_D\circ r_{\cX}\) for every vertical \(\mathbf R\)-divisor
  \(D\) on \(\cX\);
\item\label{itm:ch34-retr:ideal}
  \(\log|\mathfrak a|\le\log|\mathfrak a|\circ r_{\cX}\) for every coherent
  ideal \(\mathfrak a\) on \(\cX\) with zero locus inside \(\cX_0\);
\item\label{itm:ch34-retr:pl}
  every \(\psi\in\PL_{\mathbf R}(X)\) restricts to a real piecewise
  affine function on \(\Delta_{\cX}\), and is affine on each face when it is
  determined on \(\cX\).
\end{enumerate}
\par
\end{proposition}

\begin{proof}
Applying \Cref{lem:retraction-well-defined} and
\Cref{prop:skeleton-retraction} to the SNC \(R\)-model \(\cX_R\) produces a
continuous idempotent retraction
\(r:X_K^{\An}\to\Delta_{\cX_R}\) given by the displayed recipe.  Composing
with the Gauss section gives \(r\circ\sigma:X^{\An}\to\Delta_{\cX_R}\),
which lands in \(\sigma(X^{\An})\) by \Cref{lem:ch34-skeleton-in-gauss}; set
\(r_{\cX}:=\sigma^{-1}\circ r\circ\sigma\).  It is continuous as a
composition of continuous maps, idempotent because \(r\) is, and restricts
to the identity on \(\Delta_{\cX}\) because \(r\) does on
\(\Delta_{\cX_R}\).

Assertion \cref{itm:ch34-retr:divisor} is \Cref{lem:model-function-affine},
which gives \(\varphi_D=\varphi_D\circ r\) on \(X_K^{\An}\); composing with
\(\sigma\) and using \(\phi_D=\varphi_D\circ\sigma\) gives the statement.
Assertion \cref{itm:ch34-retr:ideal} is
\Cref{cor:retraction-ideal-monotone}, transported in the same way.
Assertion \cref{itm:ch34-retr:pl} is \Cref{prop:skeleton-model-pa} combined
with \Cref{thm:tc-pl-correspondence}: a real PL function is
\(\phi_{D'}\) for a vertical \(\mathbf R\)-divisor on some test
configuration, and passing to a common SNC domination makes it affine on
the faces of a subdivision of \(\Delta_{\cX}\).
\end{proof}

\begin{remark}
\label{rem:ch34-dual-complex-scope}
\Cref{prop:ch34-retraction} is an \emph{interpretation} of PL functions, not
their definition.  Three limitations should be kept in mind.  Existence of
SNC test configurations requires resolution of singularities, hence
characteristic zero.  Even then, a single PL function is affine on the faces
of \(\Delta_{\cX}\) only after passing to a model on which its divisor is
determined, so ``piecewise'' cannot be dropped from the name.  Each real PL
function is affine on a rational model subdivision; the failure of
\(\PL_{\mathbf R}(X)\) to be a lattice is instead that the maximum of two
such functions can have an irrational wall not realized by any rational
model subdivision.  Compare
\Cref{lem:ch34-pl-lattice}\,\cref{itm:ch34-pl:realnotlattice} and
\Cref{rem:divisorial-normalization}.
\end{remark}

\subsection{A dual complex computed: the deformation to the normal cone}

\begin{example}
\label{ex:ch34-normal-cone-segment}
Let \(X=\mathbf P^1_k\), fix a closed point \(p\in X(k)\) with local
parameter \(z\), and let
\[
 \cY:=\operatorname{Bl}_{(p,0)}\bigl(\mathbf P^1\times\mathbf A^1\bigr)
 \xrightarrow{\ \mu\ }\mathbf P^1\times\mathbf A^1 .
\]
This is the deformation to the normal cone of \Cref{ex:product-dnc-test-configurations},
and it is an SNC test configuration: \(\cY\) is a smooth surface and
\[
 \cY_0=\widetilde P+E,
 \qquad b_{\widetilde P}=b_E=1,
\]
where \(\widetilde P\) is the strict transform of \(\mathbf P^1\times\{0\}\)
and \(E\cong\mathbf P^1\) is the exceptional curve; they meet transversally
at one point \(\xi\).

\emph{The complex.}  By \Cref{lem:ch34-skeleton-in-gauss} the dual complex
\(\Delta_{\cY}\) is a segment with the two vertices
\[
 v_{\widetilde P}=v_{X,\triv},
 \qquad
 v_E=\ord_p .
\]
Indeed \(\ord_{\widetilde P}\) is trivial on \(k(X)^\times\) and takes the
value \(1\) on \(\varpi\), so it is the Gauss extension of \(v_{X,\triv}\);
and \(\ord_E(\varpi)=1\) while \(\ord_E(z)=1\), so
\(v_E=\ord_p\).  Concretely, in the chart with coordinates \((z,w)\) and
\(\varpi=zw\) one has \(E=\{z=0\}\) and \(\widetilde P=\{w=0\}\), and for
\(t\in[0,1]\) the monomial valuation with
\(\alpha_E=t\), \(\alpha_{\widetilde P}=1-t\) restricts on \(k(X)\) to
\(t\ord_p\).  Thus, in the tree coordinates of \Cref{ex:triv-p1-tree},
\begin{equation}
\label{eq:ch34-segment}
 \Delta_{\cY}=\{t\ord_p\mid 0\le t\le1\}.
\end{equation}

\emph{The retraction.}  For \(q\ne p\) the Gauss extension of \(s\ord_q\)
has center on \(\widetilde P\) only, so
\(r_{\cY}(s\ord_q)=v_{X,\triv}\) for every \(s\ge0\); and
\(r_{\cY}(v_{q,\triv})=v_{X,\triv}\) as well.  For \(s\ge1\) the center lies
on \(E\setminus\widetilde P\), so \(r_{\cY}(s\ord_p)=\ord_p\); and
\(r_{\cY}(s\ord_p)=s\ord_p\) for \(0\le s\le1\).

\emph{A vertical divisor.}  Let \(D=a\widetilde P+cE\).  By
\eqref{eq:ch34-affine-formula} and
\Cref{prop:ch34-retraction}\,\cref{itm:ch34-retr:divisor},
\begin{equation}
\label{eq:ch34-segment-function}
 \phi_D(s\ord_q)=
 \begin{cases}
 a(1-s)+cs,& q=p,\ 0\le s\le1,\\
 c,& q=p,\ s\ge1,\\
 a,& q\ne p .
 \end{cases}
\end{equation}
Since \(\cY_0=\widetilde P+E\) is numerically trivial on the fibres of
\(\pi\), replacing \(D\) by \(D-a\cY_0\) changes \(\phi_D\) by the constant
\(-a\); we therefore normalize \(a=0\) below.

\emph{Identification with Fubini--Study data.}  Take \(a=0\) and \(c=-1\).
Then \(\phi_D=\max\{-s,-1\}\) on the ray of \(p\) and \(0\) elsewhere, that
is,
\[
 \phi_D=\max\{\log|\mathfrak m_p|,-1\}
 =\phi_{\mathfrak m_p+(\varpi)},
\]
the function \(\phi_{\{p\}}\) of \Cref{ex:ch34-pl-subscheme}.  It is also a
genuine Fubini--Study function: with \(L=\cO(1)\), \(m=1\), sections
\(s_p,s_q\in H^0(X,L)\) vanishing at \(p\) and at some \(q\ne p\), and
weights \(0\) and \(-1\),
\begin{equation}
\label{eq:ch34-segment-fs}
 \max\{\log|s_p|,\ \log|s_q|-1\}=\phi_D\in\cH(\cO(1)).
\end{equation}
Indeed at \(s\ord_p\) the left side is \(\max\{-s,-1\}\); at \(s\ord_q\) it
is \(\max\{0,-s-1\}=0\); at \(s\ord_r\) with \(r\notin\{p,q\}\) it is
\(\max\{0,-1\}=0\); and at the proper-support points \(v_{r,\triv}\) it is
\(0\) for \(r\ne p\) and \(\max\{-\infty,-1\}=-1\) for \(r=p\).  In
particular the function is finite everywhere, as it must be: \(s_p\) and
\(s_q\) have no common zero.  Removing \(s_q\) from the system leaves the
generically finite function \(\log|s_p|\) of \Cref{ex:ch34-log-s-p1}, with a
pole at \(v_{p,\triv}\).
\end{example}

\begin{example}
\label{ex:ch34-two-corner}
Continue with \(X=\mathbf P^1\) and \(\cY\) as above, and let
\[
 \cY'=\operatorname{Bl}_\xi\cY\xrightarrow{\ \mu_2\ }\cY,
 \qquad \xi=\widetilde P\cap E,
\]
with exceptional curve \(G\).  Then \(\cY'\) is again a smooth surface and
\begin{equation}
\label{eq:ch34-two-corner-central-fibre}
 \cY'_0=\widetilde P'+E'+2G,
\end{equation}
because \(\ord_G(\varpi)=\ord_G(\widetilde P)+\ord_G(E)=2\).  The three
components have multiplicities \(b_{\widetilde P'}=b_{E'}=1\) and
\(b_G=2\), and \(\widetilde P'\cap E'=\varnothing\).  Hence
\(\Delta_{\cY'}\) is a segment with three vertices,
\[
 v_{\widetilde P'}=v_{X,\triv},
 \qquad
 v_G=\tfrac12\ord_p,
 \qquad
 v_{E'}=\ord_p,
\]
the middle one because \(\ord_G(z)=1\) and \(b_G=2\).  In the coordinate
\(t\) of \eqref{eq:ch34-segment}, the two faces are
\(t\in[0,\tfrac12]\) and \(t\in[\tfrac12,1]\), and on them the affine
parameters of \eqref{eq:ch34-affine-formula} are
\[
 \alpha_{\widetilde P'}=1-2t,\quad\alpha_G=t
 \qquad\text{and}\qquad
 \alpha_G=1-t,\quad\alpha_{E'}=2t-1
\]
respectively; both are affine in \(t\), so \(t\) is an affine coordinate on
each face, although the two integral affine structures differ by a factor
\(2\) in the lattice length.

\emph{A PL function with two corners.}  Take
\[
 D=cG+eE',
 \qquad
 c=-2,\quad e=-\tfrac32 .
\]
By \eqref{eq:ch34-affine-formula} the vertex values are
\(\phi_D(v_{\widetilde P'})=0\), \(\phi_D(v_G)=c/2=-1\), and
\(\phi_D(v_{E'})=e=-\tfrac32\).  Interpolating affinely on the two faces and
using \Cref{prop:ch34-retraction}\,\cref{itm:ch34-retr:divisor} beyond
\(t=1\) gives
\begin{equation}
\label{eq:ch34-two-corner-function}
 \phi_D(t\ord_p)=
 \begin{cases}
 -2t,&0\le t\le\tfrac12,\\[1mm]
 -\tfrac12-t,&\tfrac12\le t\le1,\\[1mm]
 -\tfrac32,&t\ge1,
 \end{cases}
 \qquad
 \phi_D\equiv0\ \text{off the ray of }p .
\end{equation}
The slopes are \(-2\), \(-1\), \(0\): the function is convex on the ray,
with corners exactly at the two nontrivial vertices of \(\Delta_{\cY'}\).
Equivalently
\begin{equation}
\label{eq:ch34-two-corner-max}
 \phi_D=\max\Bigl\{2\log|\mathfrak m_p|,\
 \log|\mathfrak m_p|-\tfrac12,\ -\tfrac32\Bigr\},
\end{equation}
so that \(2\phi_D=\phi_{\mathfrak a}\) for the flag ideal
\begin{equation}
\label{eq:ch34-two-corner-flag}
 \mathfrak a=\mathfrak m_p^4+\mathfrak m_p^2\varpi
 +\mathfrak m_p^2\varpi^2+\varpi^3 .
\end{equation}
Its Rees valuations, in the sense of \Cref{chap:bj-test-configurations}, are
\(\tfrac12\ord_p\) and \(\ord_p\), the two nontrivial vertices, as they must
be.

\emph{Fubini--Study data.}  Put \(L=\cO_{\mathbf P^1}(2)\) and \(m=2\), so
that \(mL=\cO(4)\).  With \(s_p,s_q\in H^0(X,\cO(1))\) as before,
\begin{equation}
\label{eq:ch34-two-corner-fs}
 \phi_D=\frac12\max\bigl\{
   \log|s_p^4|,\
   \log|s_p^2s_q^2|-1,\
   \log|s_q^4|-3\bigr\}
 \in\cH\bigl(\cO(2)\bigr).
\end{equation}
Let us verify this on each ray.  Along \(t\ord_p\) one has
\(\log|s_p|=-t\) and \(\log|s_q|=0\), so the right side is
\(\tfrac12\max\{-4t,-2t-1,-3\}\); the three affine functions of \(t\) cross
at \(t=\tfrac12\) and \(t=1\), and the maximum is \(-4t\), then
\(-2t-1\), then \(-3\), which after halving is exactly
\eqref{eq:ch34-two-corner-function}.  Along \(t\ord_q\) it is
\(\tfrac12\max\{0,-2t-1,-4t-3\}=0\); along any other ray it is
\(\tfrac12\max\{0,-1,-3\}=0\).  At \(v_{p,\triv}\) it is
\(\tfrac12\max\{-\infty,-\infty,-3\}=-\tfrac32\), finite, in accordance with
the fact that \(s_p^4,s_p^2s_q^2,s_q^4\) have no common zero.  The weights
\(0,-1,-3\) are integers, so
\(\phi_D\in\cH_{\mathbf Z}(\cO(2))\subset\cH(\cO(2))\), and the
corresponding flag ideal is \eqref{eq:ch34-two-corner-flag}, for which
\(\cO(4)\otimes\mathfrak a_j\) is
\(\cO,\cO(2),\cO(2),\cO(4)\) for \(j=0,1,2,3\), all globally generated as
\Cref{prop:ch34-fs-flag-ideal} requires.
\end{example}

Positivity of the two functions just computed is checked in
\Cref{ex:ch34-psh-check} below, once relative nefness has been defined.
\section{PL positivity}
\label{sec:ch34-pl-positivity}

Let \(\theta\in N^1(X)\).  If a test configuration \(\cX\) dominates the
trivial one and \(D\) is a vertical real Cartier divisor, denote by
\(\theta_{\cX}+D\) the corresponding numerical class relative to
\(\mathbf A^1\); here \(\theta_{\cX}\) is the pullback of \(\theta\) under
\(\cX\to X\times\mathbf A^1\to X\), and the sum is taken in
\(N^1(\cX/\mathbf A^1)\).

\begin{definition}
\label{def:positive-pl-function}
The PL function \(\phi_D\) is \(\theta\)-psh at PL level if
\begin{equation}
\label{eq:pl-relative-nefness}
 \theta_{\cX}+D\in\Nef(\cX/\mathbf A^1).
\end{equation}
\end{definition}

We write \(\PL_{\mathbf R}(X)\cap\PSH_X(\theta)\) and
\(\PL(X)\cap\PSH_X(\theta)\) for the sets of \(\theta\)-psh functions in
\(\PL_{\mathbf R}(X)\) and \(\PL(X)\).  At this point the notation is
formal; \Cref{prop:ch34-pl-consistency} will justify it by showing that it
agrees with \Cref{def:bj-psh}.

\begin{lemma}
\label{lem:ch34-pl-positivity-well-posed}
Condition \eqref{eq:pl-relative-nefness} does not depend on the choice of a
test configuration \(\cX\) and a divisor \(D\) with \(\phi_D=\phi\).
Moreover, if \eqref{eq:pl-relative-nefness} holds for some \(\theta\), then
\(\theta\in\Nef(X)\).
\end{lemma}

\begin{proof}
Let \(\mu:\cX'\to\cX\) be a morphism of test configurations.  Then
\[
 \theta_{\cX'}+\mu^*D=\mu^*(\theta_{\cX}+D).
\]
If the right-hand class is nef, so is its pullback.  Conversely, suppose
\(\mu^*(\theta_{\cX}+D)\) is nef and let \(C\subset\cX_0\) be an integral
curve.  Since \(\mu\) is proper, surjective and an isomorphism over
\(\mathbf G_m\), there is an integral curve \(C'\subset\cX'_0\) with
\(\mu(C')=C\); the projection formula gives
\((\theta_{\cX}+D)\cdot C=\deg(C'/C)^{-1}
\bigl(\mu^*(\theta_{\cX}+D)\cdot C'\bigr)\ge0\).  Hence
\eqref{eq:pl-relative-nefness} is insensitive to passing to a higher test
configuration; and two determinations of the same PL function agree after
pullback to a common domination, by
\Cref{thm:tc-pl-correspondence}\ and the injectivity established in its
proof.

For the last assertion, restrict \eqref{eq:pl-relative-nefness} to a general
fibre \(\cX_t\cong X\), \(t\ne0\), on which \(D\) is trivial: nefness is
preserved by restriction to a subvariety, so
\(\theta\in\Nef(X)\).
\end{proof}

\begin{lemma}
\label{lem:ch34-pl-limit-classes}
Let \(\theta_i\to\theta\) in \(N^1(X)\), let
\(\phi,\psi\in\PL_{\mathbf R}(X)\), and let \(c_i\to0\) in \(\mathbf R\).
If \(\phi+c_i\psi\) is \(\theta_i\)-psh at PL level for every \(i\), then
\(\phi\) is \(\theta\)-psh at PL level.
\end{lemma}

\begin{proof}
By \Cref{thm:tc-pl-correspondence} and
\Cref{lem:ch34-pl-positivity-well-posed} we may determine \(\phi=\phi_D\)
and \(\psi=\phi_{D'}\) on one test configuration \(\cX\).  The hypothesis
says \((\theta_i)_{\cX}+D+c_iD'\in\Nef(\cX/\mathbf A^1)\).  The relative nef
cone is a closed convex cone in the finite-dimensional space
\(N^1(\cX/\mathbf A^1)\), and
\((\theta_i)_{\cX}+D+c_iD'\to\theta_{\cX}+D\); hence the limit is nef.
\end{proof}

\begin{proposition}
\label{prop:ch34-pl-operations}
Let \(\theta,\theta'\in N^1(X)\) and let
\(\phi,\phi'\in\PL_{\mathbf R}(X)\) be \(\theta\)-psh and \(\theta'\)-psh at
PL level.  Then:
\begin{enumerate}
\item\label{itm:ch34-plop:sum}
  \(\phi+\phi'\) is \((\theta+\theta')\)-psh, and \(t\phi\) is
  \(t\theta\)-psh for \(t\in\mathbf R_{>0}\);
\item\label{itm:ch34-plop:const}
  \(\phi+c\) is \(\theta\)-psh for every \(c\in\mathbf R\), and
  \(t\cdot\phi\) is \(\theta\)-psh for every \(t\in\mathbf Q_{>0}\), where
  the action is that of \eqref{eq:function-scaling-action};
\item\label{itm:ch34-plop:max}
  if \(\theta=\theta'\) and \(\phi,\phi'\in\PL(X)\), then
  \(\max\{\phi,\phi'\}\in\PL(X)\) is \(\theta\)-psh;
\item\label{itm:ch34-plop:decreasing}
  \(\phi\) is decreasing for the order \eqref{eq:semivaluation-order}, and
  therefore satisfies the maximum principle
  \(\sup_{Y^{\An}}\phi=\phi(v_{Y,\triv})\) for every irreducible
  \(Y\subset X\);
\item\label{itm:ch34-plop:pullback}
  for every morphism \(f:Y\to X\) from a projective variety, \(f^*\phi\) is
  \(f^*\theta\)-psh at PL level.
\end{enumerate}
Consequently \(\PL_{\mathbf R}(X)\cap\PSH_X(\theta)\) is a convex subset of
\(\PL_{\mathbf R}(X)\), closed for the topology of uniform convergence.
\end{proposition}

\begin{proof}
Determine \(\phi=\phi_D\) and \(\phi'=\phi_{D'}\) on a common test
configuration \(\cX\), so that \(\theta_{\cX}+D\) and
\(\theta'_{\cX}+D'\) are relatively nef.

\emph{\cref{itm:ch34-plop:sum}.}  Then
\((\theta+\theta')_{\cX}+D+D'=(\theta_{\cX}+D)+(\theta'_{\cX}+D')\) and
\((t\theta)_{\cX}+tD=t(\theta_{\cX}+D)\) are nef, the relative nef cone
being a convex cone.

\emph{\cref{itm:ch34-plop:const}.}  We have \(\phi+c=\phi_{D+c\cX_0}\), and
\(\cX_0=\pi^*(0)\) vanishes in \(N^1(\cX/\mathbf A^1)\); so
\(\theta_{\cX}+D+c\cX_0=\theta_{\cX}+D\) is nef.  For the scaling action it
suffices to treat \(t=d\in\mathbf Z_{>0}\) and \(t=1/d\).  By
\Cref{prop:ch34-fs-operations}\,\cref{itm:ch34-ops:scaling} and
\eqref{eq:ch34-flag-weight-form}, \(d\cdot\phi_D=\phi_{D_d}\), where
\(D_d\) is the pullback of \(D\) under the degree-\(d\) base change
\(\cX_d\to\cX\), \(\varpi\mapsto\varpi^d\); this is
\cite[Lemma~2.34]{BJ22}.  A class on \(\cX\) is relatively nef if and only
if its pullback to \(\cX_d\) is, by the argument of
\Cref{lem:ch34-pl-positivity-well-posed} applied to the finite surjective
morphism \(\cX_d\to\cX\).  Hence \(d\cdot\phi\) is \(\theta\)-psh, and
applying the same with \(d^{-1}\) gives the general rational case.

\emph{\cref{itm:ch34-plop:max}.}  Assume \(\theta=\theta'\) and
\(\phi,\phi'\in\PL(X)\).  Fix an ample \(\mathbf Q\)-line bundle \(L\) with
\(c_1(L)-\theta\) ample and as small as we wish.  Since
\((c_1(L))_{\cX}+D=(\theta_{\cX}+D)+(c_1(L)-\theta)_{\cX}\) is a sum of a
nef class and the pullback of a nef class, \(\phi\) and \(\phi'\) are
\(L\)-psh at PL level.  Passing to a higher test configuration, we may
assume by \cite[Lemma~2.33]{BJ22} that \(\cX\) carries a relatively ample
test configuration \(\cL\) for \(L\); write \(\phi_{\cL}\in\PL(X)\) for its
function.  For \(\varepsilon\in\mathbf Q_{>0}\) the class
\[
 (c_1(L))_{\cX}+D+\varepsilon\bigl(\cL-(c_1(L))_{\cX}\bigr)
\]
is relatively ample, hence relatively semiample, so by
\Cref{thm:tc-pl-correspondence} the function
\(\phi+\varepsilon\phi_{\cL}\) lies in \(\cH\bigl((1+\varepsilon)L\bigr)\);
the same holds for \(\phi'+\varepsilon\phi_{\cL}\).  By
\Cref{lem:gf-fs-operations}\,(2),
\[
 \max\{\phi,\phi'\}+\varepsilon\phi_{\cL}
 =\max\{\phi+\varepsilon\phi_{\cL},\ \phi'+\varepsilon\phi_{\cL}\}
 \in\cH\bigl((1+\varepsilon)L\bigr),
\]
where the first equality is \Cref{lem:extended-real-safe}\,(1) applied
pointwise, all values being finite.  A Fubini--Study function is
\(L\)-psh at PL level by \Cref{thm:tc-pl-correspondence}, semiample
implying nef; so \(\max\{\phi,\phi'\}+\varepsilon\phi_{\cL}\) is
\((1+\varepsilon)c_1(L)\)-psh.  Letting \(\varepsilon\downarrow0\) and
applying \Cref{lem:ch34-pl-limit-classes} with
\(\psi=\phi_{\cL}\) shows that \(\max\{\phi,\phi'\}\) is \(L\)-psh; letting
\(c_1(L)\downarrow\theta\) and applying
\Cref{lem:ch34-pl-limit-classes} again gives the claim.

\emph{\cref{itm:ch34-plop:decreasing}.}  By
\Cref{lem:ch34-pl-fs-approximation} below, \(\phi\) is a uniform limit of
Fubini--Study functions, each of which is decreasing by
\Cref{prop:ch34-generic-finiteness}\,\cref{itm:ch34-gf:decreasing}; a uniform
limit of decreasing functions is decreasing, and the maximum principle
follows as in the proof of that proposition.

\emph{\cref{itm:ch34-plop:pullback}.}  The equivariant rational map
\(Y\times\mathbf G_m\dashrightarrow\cX\) induced by \(f\) has an equivariant
resolution of indeterminacy, which is a test configuration \(\cY\) for \(Y\)
dominating \(Y\times\mathbf A^1\) together with a morphism
\(g:\cY\to\cX\).  Then \(f^*\phi=\phi_{g^*D}\) and
\((f^*\theta)_{\cY}+g^*D=g^*(\theta_{\cX}+D)\) is nef.

The final sentence follows from \cref{itm:ch34-plop:sum} and
\cref{itm:ch34-plop:const} for convexity, and from
\Cref{lem:ch34-pl-limit-classes} for closedness: if
\(\phi_j\to\phi\) uniformly with all \(\phi_j\) \(\theta\)-psh, determine
them on a common test configuration and use that
\(D_j\to D\) in the finite-dimensional space
\(\operatorname{VCar}(\cX)_{\mathbf R}\), by injectivity of
\(D\mapsto\phi_D\) and the equivalence of norms there.
\end{proof}

\begin{lemma}
\label{lem:ch34-pl-fs-approximation}
Let \(\theta\in N^1(X)\) and let \(\phi\in\PL_{\mathbf R}(X)\) be
\(\theta\)-psh at PL level.  Then there exist \(\mathbf Q\)-line bundles
\(L_j\) with \(c_1(L_j)\to\theta\) and \(\phi_j\in\cH(L_j)\) with
\(\phi_j\to\phi\) uniformly on \(X^{\An}\).
\end{lemma}

\begin{proof}
Determine \(\phi=\phi_D\) on a test configuration \(\cX\) dominating
\(X\times\mathbf A^1\), with \(D\in\operatorname{VCar}(\cX)_{\mathbf R}\) and
\(\theta_{\cX}+D\) relatively nef.  The relatively ample cone is open in
\(N^1(\cX/\mathbf A^1)\) and its closure is the relatively nef cone, so we
may choose \(\mathbf Q\)-line bundles \(L_j\) on \(X\) and
\(D_j\in\operatorname{VCar}(\cX)_{\mathbf Q}\) with
\(c_1(L_j)\to\theta\) in \(N^1(X)\), \(D_j\to D\) in
\(\operatorname{VCar}(\cX)_{\mathbf R}\), and
\((c_1(L_j))_{\cX}+D_j\) relatively ample.  Then
\(\cL_j:=(L_j)_{\cX}+D_j\) is a semiample test configuration for \(L_j\),
so \(\phi_j:=\phi_{D_j}\in\cH(L_j)\) by \Cref{thm:tc-pl-correspondence}.
Finally \(D\mapsto\phi_D\) is a linear map from the finite-dimensional
space \(\operatorname{VCar}(\cX)_{\mathbf R}\) to the Banach space
\(C^0(X^{\An})\), hence bounded; so \(D_j\to D\) implies
\(\phi_j\to\phi\) uniformly.  This is \cite[Lemma~3.7]{BJ22}.
\end{proof}

For an ample class it is convenient to have a distinguished cofinal family
of PL psh functions.

\begin{definition}
\label{def:ch34-hdom}
For \(\omega\in\Amp(X)\), let \(\cH_{\mathrm{dom}}(\omega)\) be the set of
\(\phi_D\) with \(D\in\operatorname{VCar}(\cX)_{\mathbf Q}\) for some test
configuration \(\cX\) dominating \(X\times\mathbf A^1\) and
\(\omega_{\cX}+D\) relatively ample.  If \(\omega=c_1(L)\) with \(L\) an
ample \(\mathbf Q\)-line bundle we write
\(\cH_{\mathrm{dom}}(L):=\cH_{\mathrm{dom}}(c_1(L))\subset\cH(L)\).
\end{definition}

By \Cref{cor:ch34-pl-plus-union} and the openness of the relatively ample
cone, \(\mathbf Q_{>0}\cH_{\mathrm{dom}}(\omega)=\PL^+(X)\), so
\(\cH_{\mathrm{dom}}(\omega)\) spans \(\PL(X)\).  Moreover every
\(\phi\in\PL_{\mathbf R}(X)\cap\PSH_X(\omega)\) is a uniform limit of
elements of \(\cH_{\mathrm{dom}}(\omega)\); this is
\cite[Proposition~3.10]{BJ22}, whose proof is the argument of
\Cref{lem:ch34-pl-fs-approximation} with the class \(\omega\) held fixed and
the perturbation absorbed into a relatively ample auxiliary divisor.

\begin{warning}{Nef is not semiample}
The condition \eqref{eq:pl-relative-nefness} is genuinely weaker than
relative semiampleness.  For rational data, relative semiampleness implies
\eqref{eq:pl-relative-nefness}, and the corresponding function lies in
\(\cH(L)\) by \Cref{thm:tc-pl-correspondence}; the converse can require
approximation and must not be used without hypotheses.  Concretely,
\(\cH(L)\subset\PL(X)\cap\PSH_X(c_1(L))\) and the inclusion is strict in
general, which is why \Cref{lem:ch34-pl-fs-approximation} produces a family
of \emph{varying} classes \(L_j\) rather than a single one.
\end{warning}

\subsection{The two examples, checked}

\begin{example}
\label{ex:ch34-psh-check}
We check \eqref{eq:pl-relative-nefness} for the functions computed in
\Cref{ex:ch34-normal-cone-segment,ex:ch34-two-corner}.  Throughout
\(X=\mathbf P^1\) and \(\theta=c_1(\cO(d))\), so that \(\theta\) is
determined by the integer \(d=\deg\theta\); on a test configuration \(\cY\)
the class \(\theta_{\cY}\) is the pullback of \(\cO(d)\) under
\(\cY\to\mathbf P^1\).

\emph{The two-vertex case.}  On the smooth surface
\(\cY=\operatorname{Bl}_{(p,0)}(\mathbf P^1\times\mathbf A^1)\) the relevant
intersection numbers are
\[
 E^2=-1,\qquad \widetilde P^2=-1,\qquad E\cdot\widetilde P=1,
 \qquad
 \theta_{\cY}\cdot\widetilde P=d,\qquad\theta_{\cY}\cdot E=0 ,
\]
the last because \(E\) is contracted by \(\cY\to\mathbf P^1\), and
\(\widetilde P^2=-1\) because \(\mathbf P^1\times\{0\}\) is a fibre, of
self-intersection \(0\), through the blown-up point.  With \(D=cE\), the
curves contained in fibres of \(\cY\to\mathbf A^1\) are \(\widetilde P\),
\(E\), and the general fibres \(\cY_t\cong\mathbf P^1\); hence
\eqref{eq:pl-relative-nefness} amounts to
\[
 (\theta_{\cY}+D)\cdot\widetilde P=d+c\ \ge0,
 \qquad
 (\theta_{\cY}+D)\cdot E=-c\ \ge0,
 \qquad
 (\theta_{\cY}+D)\cdot\cY_t=d\ \ge0 .
\]
So \(\phi_D\) is \(\theta\)-psh at PL level exactly when
\begin{equation}
\label{eq:ch34-segment-psh-range}
 -d\le c\le0 .
\end{equation}
With \(d=1\) and \(c=-1\) both inequalities are equalities, and
\eqref{eq:ch34-segment-fs} exhibits the corresponding function as an element
of \(\cH(\cO(1))\); this is the extreme case, in which the slope of
\(\phi_D\) at the root of the ray equals \(-\deg\theta\).

\emph{The three-vertex case.}  On \(\cY'=\operatorname{Bl}_\xi\cY\), writing
\(\widetilde P'=\mu_2^*\widetilde P-G\) and \(E'=\mu_2^*E-G\), one computes
\[
 G^2=-1,\quad \widetilde P'^2=E'^2=-2,\quad
 \widetilde P'\cdot G=E'\cdot G=1,\quad
 \widetilde P'\cdot E'=0,
\]
and \(\theta_{\cY'}\cdot\widetilde P'=d\),
\(\theta_{\cY'}\cdot G=\theta_{\cY'}\cdot E'=0\).  As a check,
\(\cY'_0=\widetilde P'+E'+2G\) has zero intersection with each of the three
components:
\(\cY'_0\cdot G=1+1-2=0\),
\(\cY'_0\cdot E'=0-2+2=0\),
\(\cY'_0\cdot\widetilde P'=-2+0+2=0\).
With \(D=cG+eE'\), condition \eqref{eq:pl-relative-nefness} becomes
\begin{equation}
\label{eq:ch34-two-corner-nef}
 \begin{aligned}
  (\theta_{\cY'}+D)\cdot\widetilde P'&=d+c\ge0,\\
  (\theta_{\cY'}+D)\cdot G&=-c+e\ge0,\\
  (\theta_{\cY'}+D)\cdot E'&=c-2e\ge0.
 \end{aligned}
\end{equation}
For the data \(c=-2\), \(e=-\tfrac32\) of \Cref{ex:ch34-two-corner} the
three numbers are
\[
 d-2,\qquad \tfrac12,\qquad 1 .
\]
Hence \(\phi_D\) is \(\theta\)-psh at PL level exactly for
\(d\ge2\), and \(d=2\) is the borderline case.  This is consistent with
\eqref{eq:ch34-two-corner-fs}, which exhibits \(\phi_D\) as an element of
\(\cH(\cO(2))\), hence as a \(c_1(\cO(2))\)-psh PL function by
\Cref{thm:tc-pl-correspondence}.

\emph{Reading the conditions on the graph.}  Recall from
\eqref{eq:ch34-two-corner-function} that the slopes of
\(t\mapsto\phi_D(t\ord_p)\) are \(-2\), \(-1\), \(0\) on the three
intervals \([0,\tfrac12]\), \([\tfrac12,1]\), \([1,\infty)\).  In terms of
those slopes, the three inequalities \eqref{eq:ch34-two-corner-nef} read
\[
 \deg\theta+\phi_D'(0^+)\ge0,
 \qquad
 \phi_D'(\tfrac12^-)\le\phi_D'(\tfrac12^+),
 \qquad
 \phi_D'(1^-)\le\phi_D'(1^+),
\]
that is: one global inequality involving the degree of the class, and one
convexity inequality at each corner.  The same reformulation applies in the
two-vertex case, where \eqref{eq:ch34-segment-psh-range} says
\(d+c\ge0\) and \(c\le0\).  These are instances of the general criterion for
smooth projective curves, \cite[Example~4.13]{BJ22}: a function
\(\phi:X^{\An}\to\mathbf R\cup\{-\infty\}\) is \(\theta\)-psh if and only if
\(t\mapsto\phi(t\ord_q)\) is convex on \([0,+\infty)\) for every closed
point \(q\), and
\(\deg\theta+\sum_{q}\phi'_q(0^+)\ge0\), the sum having only countably many
nonzero terms.
\end{example}

\section{The class \texorpdfstring{\(\PSH_X(\theta)\)}{PSH(theta)}}
\label{sec:ch34-psh}

Everything so far has been finitary: a Fubini--Study function is built from
finitely many sections, a PL function from a single test configuration.  The
class of psh functions is obtained by closing up under decreasing limits,
and it is the smallest class for which this is done.

\begin{definition}
\label{def:bj-psh}
Let \(\theta\in N^1(X)\).  A function
\(\phi:X^{\An}\to\mathbf R\cup\{-\infty\}\) is \emph{\(\theta\)-psh} if
\begin{enumerate}
\item \(\phi\) is upper semicontinuous;
\item \(\phi\) is generically finite, that is, \(\phi\not\equiv-\infty\) on
      any irreducible component of \(X\); and
\item there are \(\mathbf Q\)-line bundles \(L_i\) and a decreasing net
      \(\phi_i\in\cH^{\mathrm{gf}}_{\mathbf R}(L_i)\) with
      \(\lim_i c_1(L_i)=\theta\) in \(N^1(X)\) and \(\phi=\lim_i\phi_i\)
      pointwise on \(X^{\An}\).
\end{enumerate}
The set of \(\theta\)-psh functions is denoted \(\PSH_X(\theta)\).
\end{definition}

Three points of the definition deserve comment before anything is proved
with it.

First, condition~(3) uses a \emph{net}, not a sequence.  This is not
pedantry: \(X^{\An}\) is compact but not metrizable in general, the index
set of the natural approximating family is the set of test configurations
ordered by domination, and that set is directed but not countable.  The
places where a sequence \emph{may} legitimately be extracted are recorded in
\Cref{chap:bj-linear-growth}; none of them is available here.

Second, the limit in~(3) is a decreasing limit, so it exists in
\(\Rbar\) at every point without hypothesis, by ledger entry~(1), and the
limit function is automatically upper semicontinuous as an infimum of
continuous functions.  Condition~(1) is therefore redundant given~(3); it is
retained because it is the condition one verifies in practice when a
candidate \(\phi\) is presented by other means.

Third, condition~(2) is a genuine nontriviality condition.  It excludes the
function \(\equiv-\infty\), and on a reducible \(X\) it is imposed
componentwise: being finite on only some components is not enough.

Fourth, and most importantly, the classes are allowed to \emph{vary}: the
approximants live in \(\cH^{\mathrm{gf}}_{\mathbf R}(L_i)\) for line bundles
whose classes merely converge to \(\theta\).  This is not a technical
convenience: it is part of the closure construction of
\cite[Definition~4.1]{BJ22} and allows arbitrary numerical limit classes.
No claim is made that fixed-class approximation would force nefness; GF
Fubini--Study functions can occur in non-nef effective classes.  For \(E\)
an effective divisor, \(\log|s_E|\) is \([E]\)-psh, and \([E]\) is only
pseudoeffective in general.  The precise
consequence of the varying-class convention is
\Cref{lem:ch34-psh-nonempty}, which separates the pseudoeffective threshold
for \(\PSH_X(\theta)\) from the nef threshold for its \emph{bounded}
members.

\begin{lemma}
\label{lem:ch34-psh-nonempty}
Let \(\theta\in N^1(X)\).  Then
\begin{equation}
\label{eq:psh-existence-cones}
 \begin{aligned}
 \PSH_X(\theta)\ne\emptyset
 &\implies \theta\in\operatorname{Psef}(X),\\
 \PSH_X(\theta)\ \text{contains a bounded function}
 &\implies \theta\in\Nef(X),
 \end{aligned}
\end{equation}
and conversely the constants are \(\theta\)-psh whenever
\(\theta\in\Nef(X)\).  In particular
\(\mathbf R\subset\PSH_X(\theta)\) if and only if \(\theta\in\Nef(X)\).
\end{lemma}

\begin{proof}
Suppose \(\PSH_X(\theta)\ne\emptyset\) and let \((\phi_i)\) be as in
\Cref{def:bj-psh}(3).  Each \(\cH^{\mathrm{gf}}_{\mathbf R}(L_i)\) is
nonempty, so \(L_i\) admits a nonzero section after clearing denominators
and is therefore effective; hence \(c_1(L_i)\in\operatorname{Psef}(X)\).  The
pseudoeffective cone is closed, and \(c_1(L_i)\to\theta\), so
\(\theta\in\operatorname{Psef}(X)\).

Suppose now that some \(\phi\in\PSH_X(\theta)\) is bounded.  Every member
\(\phi_i\) of a defining decreasing net satisfies \(\phi_i\geq\phi\), so it
is finite valued, hence
\(\phi_i\in\cH_{\mathbf R}(L_i)\) and the defining sections have no common
zero; consequently \(L_i\) is semiample, so \(c_1(L_i)\in\Nef(X)\).  The nef
cone is closed, whence \(\theta\in\Nef(X)\).

Conversely let \(\theta\in\Nef(X)\) and \(c\in\mathbf R\).  Choose
\(\mathbf Q\)-line bundles \(L_i\) with \(c_1(L_i)\to\theta\) and \(L_i\)
ample; for each \(i\) the base-point-free sections of a large multiple
present the constant \(c\) as an element of \(\cH_{\mathbf R}(L_i)\) after
translating the weights, and the net may be arranged to decrease to \(c\).
Hence \(c\in\PSH_X(\theta)\).  Since constants are bounded, the second
implication of \eqref{eq:psh-existence-cones} shows the converse fails when
\(\theta\) is not nef.
\end{proof}

Because constants are always available when the class is nef, and because
adding a constant preserves everything in sight, it is convenient to fix a
normalization once and for all.  We set
\begin{equation}
\label{eq:normalized-psh-space}
 \PSH_X^{\sup}(\theta)
 :=\Bigl\{\phi\in\PSH_X(\theta)\ :\ \sup_{X^{\An}}\phi=0\Bigr\},
 \qquad
 \PSH_X(\theta)=\PSH_X^{\sup}(\theta)+\mathbf R ,
\end{equation}
the second identity holding because \(\sup_{X^{\An}}\phi\) is finite for
every \(\theta\)-psh \(\phi\): it is \(>-\infty\) by generic finiteness,
\Cref{def:bj-psh}(2), and \(<+\infty\) because \(\phi\) is dominated by any
member of an approximating net, and a generically finite Fubini--Study
function is upper semicontinuous on the compact space \(X^{\An}\) and hence
bounded above.  Thus \(\PSH_X^{\sup}(\theta)\) is nonempty exactly when
\(\PSH_X(\theta)\) is nonempty: any psh function can be shifted by its
finite supremum.  What requires \(\theta\in\Nef(X)\) is the presence of a
bounded member, equivalently of the constants, by
\Cref{lem:ch34-psh-nonempty}.

\begin{proposition}
\label{prop:ch34-pl-consistency}
Let \(\phi\in\PL_{\mathbf R}(X)\).  Then \(\phi\in\PSH_X(\theta)\) in the
sense of \Cref{def:bj-psh} if and only if \(\phi\) is \(\theta\)-psh at PL
level in the sense of \Cref{def:positive-pl-function}.
\end{proposition}

\begin{proof}
This is exactly \cite[Lemma~4.3]{BJ22}.  Its proof there derives both
directions from \cite[Lemmas~3.5 and~3.7 and Corollary~3.22]{BJ22}: PL
positivity is transported between a fixed model and nearby rational
classes, and the PL-positive cone is closed in the finite-dimensional
direct-limit topology.  No envelope theorem, and in particular no result
from \S5 of \cite{BJ22}, enters this comparison.
\end{proof}

\begin{remark}
The reader should notice what \Cref{prop:ch34-pl-consistency} does
\emph{not} say.  It compares two notions on the set of PL functions.  It
does not say that a decreasing limit of PL \(\theta\)-psh functions is PL,
which is false already on \(\mathbf P^1\), nor that every \(\theta\)-psh
function is a decreasing limit of a \emph{sequence}, which is exactly the
point at which nets cannot be dispensed with.
\end{remark}

\section{Basic closure properties}
\label{sec:ch34-closure}

\begin{warning}{Primary-source import; Rethlas gate}
The next theorem is an exact source-aligned import from
\cite[Definition~4.1, Theorems~4.5 and~4.7, Lemmas~4.6 and~4.9,
Corollary~4.17, Theorem~4.22, and Corollary~4.23]{BJ22}.  The present audit
found no exact-current Rethlas receipt for this repaired block.  Accordingly
none of its clauses is presented as newly machine-verified mathematics, and
no extension beyond the cited clauses may be inferred without a separate
source proof or an exact-current receipt.
\end{warning}

\begin{theorem}
\label{thm:bj-basic-closure}
Let \(\theta,\theta'\in N^1(X)\) be arbitrary numerical classes.  No
positivity hypothesis is needed: by \Cref{lem:ch34-psh-nonempty} the
assertions below are vacuous unless the relevant classes are
pseudoeffective, and nefness is required only where a \emph{bounded} member
is asserted to exist, which is nowhere in this theorem.
\begin{enumerate}
\item If \(\phi\in\PSH_X(\theta)\) and \(c\in\mathbf R\), then
      \(\phi+c\in\PSH_X(\theta)\).
\item If \(\phi,\psi\in\PSH_X(\theta)\), then
      \(\max\{\phi,\psi\}\in\PSH_X(\theta)\).
\item If \(\phi\in\PSH_X(\theta)\) and
      \(\psi\in\PSH_X(\theta')\), then
      \(\phi+\psi\in\PSH_X(\theta+\theta')\), and
      \(a\phi\in\PSH_X(a\theta)\) for every \(a>0\).  Thus strictly
      positive linear combinations add the corresponding classes.  An
      endpoint with coefficient zero means that the corresponding summand
      is omitted; it is never interpreted as \(0\cdot(-\infty)\).
\item If \((\psi_j)_{j\in J}\) is a decreasing net in \(\PSH_X(\theta)\) and
      \(\psi:=\lim_j\psi_j\) satisfies \(\psi\not\equiv-\infty\) on any
      irreducible component, then \(\psi\in\PSH_X(\theta)\).
\item \emph{(Varying classes.)}  More generally, if \(\psi_j\in\PSH_X(\theta_j)\)
      with \(\theta_j\to\theta\) in \(N^1(X)\), the net \((\psi_j)\) is
      decreasing, and \(\psi:=\lim_j\psi_j\) is generically finite, then
      \(\psi\in\PSH_X(\theta)\).  Consequently, for a fixed function
      \(\psi\) the set \(\{\theta:\psi\in\PSH_X(\theta)\}\) is closed.
\item If \(f:Y\to X\) is a morphism of projective varieties and
      \(\phi\in\PSH_X(\theta)\), then either
      \(f^*\phi\in\PSH_Y(f^*\theta)\), or \(f^*\phi\equiv-\infty\) on an
      irreducible component of \(Y\).  In particular, restriction to every
      irreducible component \(X_\alpha\) belongs to
      \(\PSH_{X_\alpha}(\theta|_{X_\alpha})\).
\item\label{itm:bj-closure:determination} Every \(\phi\in\PSH_X(\theta)\) is finite on
      \(X^{\lin}\), hence on \(X^{\divis}\).  Moreover, for every upper
      semicontinuous \(\eta:X^{\An}\to\mathbf R\cup\{-\infty\}\),
      \[
       \phi\leq\eta\ \text{on }X^{\divis}
       \quad\Longleftrightarrow\quad
       \phi\leq\eta\ \text{on }X^{\An}.
      \]
      In particular, psh functions are determined by their divisorial
      values.
\end{enumerate}
\end{theorem}

\begin{proof}
(1) If \((\phi_i)\) approximates \(\phi\) as in \Cref{def:bj-psh}(3), then
\((\phi_i+c)\) is a decreasing net with the same classes, and
\(\phi_i+c\in\cH^{\mathrm{gf}}_{\mathbf R}(L_i)\): in a presentation of
denominator \(m_i\), replacing every weight \(\lambda_j\) by
\(\lambda_j+m_i c\) in \eqref{eq:gf-fs-presentation} adds \(c\) to the
function.  No positivity of \(\theta\) is used.

(2) The point that the two approximating families must have the same class
at each stage is essential.  The simultaneous approximation lemma
\cite[Lemma~4.9]{BJ22} gives a net of \(\mathbf Q\)-line bundles \(L_i\) and
functions
\(\phi_i,\psi_i\in\cH^{\mathrm{gf}}_{\mathbf R}(L_i)\) such that
\(c_1(L_i)\to\theta\), \(\phi_i\searrow\phi\), and
\(\psi_i\searrow\psi\).  By
\Cref{lem:gf-fs-operations}\,(2),
\(\max\{\phi_i,\psi_i\}\in
\cH^{\mathrm{gf}}_{\mathbf R}(L_i)\), and these functions decrease
pointwise to \(\max\{\phi,\psi\}\).  The limit is generically finite
because \(\phi\) is, so it is \(\theta\)-psh.

(3) Choose approximating nets \((\phi_i)_{i\in I}\) and
\((\psi_{i'})_{i'\in I'}\).  On the directed product \(I\times I'\), the
net \(\phi_i+\psi_{i'}\) decreases to \(\phi+\psi\), and
\Cref{lem:gf-fs-operations}\,(1) places its terms in class
\(L_i+L'_{i'}\), whose Chern classes converge to \(\theta+\theta'\).
The limit is generically finite since both summands are finite at every
trivial valuation.  This proves the sum assertion.  Positive rational
scalars follow from
\Cref{prop:ch34-fs-operations}\,\cref{itm:ch34-ops:rational}.  For
\(a\in\mathbf R_{>0}\), subtract the finite number \(\sup\phi\) so that
\(\phi\leq0\), choose rational \(a_m\uparrow a\), and apply~(5) to the
decreasing sequence \(a_m\phi\in\PSH_X(a_m\theta)\); then restore the
finite constant \(a\sup\phi\) using~(1).  All scalar multiplications in
this argument have a strictly positive scalar.  Together with the sum
assertion this proves the stated positive-linear-combination rule.  This is
\cite[Theorem~4.7\,(i)--(ii)]{BJ22}.

(5) This is precisely property \emph{(PSH2)} of
\cite[Theorem~4.5]{BJ22}.  We spell out the diagonal issue that its proof
resolves.  Apply \cite[Lemma~4.6]{BJ22} with
\[
 \begin{aligned}
  K&=X^{\An},& S&=\Pic(X)_{\mathbf Q},
    &F_L&=\cH^{\mathrm{gf}}_{\mathbf R}(L),\\
  T&=N^1(X),& \rho(L)&=c_1(L).
 \end{aligned}
\]
Here \(K\) is compact, \(T\) is finite dimensional and hence metrizable,
and every \(F_L\) is stable under addition of a real constant by
\Cref{lem:gf-fs-operations}\,(3).  The lemma constructs the diagonal on
triples \((j,i,\varepsilon)\), ordered simultaneously by the outer index,
the pointwise order of the shifted approximants, and the class error.  That
preorder is directed by the compactness argument in the lemma and produces
one decreasing GF net whose classes converge to \(\theta\) and whose
pointwise limit is \(\psi\).  This is the required condition in
\Cref{def:bj-psh}(3).  In particular, merely ordering the disjoint union of
the inner index sets does \emph{not} prove directedness.  The closedness
assertion is the special case in which \(\psi_j=\psi\) for all \(j\).

(4) is the fixed-class special case of~(5).

(6) Write \(\phi_i\searrow\phi\) as in \Cref{def:bj-psh}.  Pullback of a
GF presentation is again GF unless all its sections vanish identically on
some component of the source, by \Cref{lem:gf-fs-operations}\,(4).  Thus,
if \(f^*\phi\) is generically finite, every \(f^*\phi_i\) is generically
finite, \(f^*\phi_i\searrow f^*\phi\), and
\(c_1(f^*L_i)\to f^*\theta\); hence
\(f^*\phi\in\PSH_Y(f^*\theta)\).  Otherwise the alternative in the
statement holds.  Restriction to \(X_\alpha\) is generically finite by
definition.  This is \cite[Theorem~4.7\,(iii)--(iv)]{BJ22}.

(7) Finiteness on \(X^{\lin}\) is
\cite[Corollary~4.17]{BJ22}.  The comparison statement is
\cite[Theorem~4.22]{BJ22}, and divisorial determination is its
\cite[Corollary~4.23]{BJ22}.
\end{proof}

\begin{warning}{Suprema are not yet available}
\Cref{thm:bj-basic-closure} contains no statement about \emph{increasing}
families, and this omission is the central open point of the trivially
valued theory as it stands at this stage of the book.  Given a nonempty
family \((\phi_\alpha)\subset\PSH_X(\theta)\) that is bounded above, the
regularized supremum \((\sup_\alpha\phi_\alpha)^*\) need not be visibly
\(\theta\)-psh: the definition supplies decreasing nets, and there is no
formal manoeuvre converting an increasing family into one.  That
\((\sup_\alpha\phi_\alpha)^*\) is nevertheless \(\theta\)-psh, for \(X\)
integral and unibranch, is the envelope theorem, stated and proved in
\Cref{chap:bj-envelope-statement}.  Until then, no argument in this book
uses closure under suprema, and the reader should treat every appeal to an
upper envelope before that chapter as a forward reference to be discharged
there.
\end{warning}

\section{Regularization and compactness}
\label{sec:ch34-compactness}

\Cref{def:bj-psh} approximates by PL functions.  For many purposes one wants
the approximants to be Fubini--Study, which is a genuinely stronger
statement because it produces sections rather than divisors.

\begin{theorem}
\label{thm:ch34-fs-regularization}
Let \(\theta\in N^1(X)\) and let \(\phi\in\PSH_X(\theta)\).  Then there are
\(\mathbf Q\)-line bundles \(L_i\) and a decreasing net
\(\phi_i\in\cH^{\mathrm{gf}}_{\mathbf Q}(L_i)\) such that
\begin{equation}
\label{eq:bj-psh-regularization}
 \begin{aligned}
  \phi&=\lim_i\phi_i &&\text{pointwise on }X^{\An},
       &\phi_i&\downarrow\phi,\\
  c_1(L_i)-\theta&\in\Amp(X),
       &&&c_1(L_i)&\longrightarrow\theta .
 \end{aligned}
\end{equation}
If \(\theta\) is nef, the approximants may be chosen in \(\cH(L_i)\) with
each \(L_i\) ample.  If \(\theta=\omega\) is ample, they may be chosen in
\(\cH_{\mathrm{dom}}(\omega)\).  Conversely, every decreasing limit in the
first display that is generically finite is \(\theta\)-psh.
\end{theorem}

\begin{proof}
The three direct assertions are precisely
\cite[Theorem~4.15\,(i)--(iii)]{BJ22}.  That proof combines the rational PL
approximation of \cite[Lemma~3.7]{BJ22}, the ample domination construction
of \cite[Proposition~3.10]{BJ22}, and the diagonal-net lemma
\cite[Lemma~4.6]{BJ22}.  The converse is the defining decreasing-closure
property \cite[Definition~4.1 and Theorem~4.5]{BJ22}.  No envelope theorem
is used; \cite[Theorem~5.20]{BJ22} concerns the envelope property and is not
a regularization result.
\end{proof}

The class drift in \eqref{eq:bj-psh-regularization} is essential: the
approximants generally live in nearby ample rational classes, not in one
fixed line-bundle class.

\begin{theorem}
\label{thm:ch34-psh-compactness}
Let \(\theta\in N^1(X)\).  Every \(\theta\)-psh function is finite on
\(X^{\divis}\) and determined by its restriction there.  Consequently, the
topology of pointwise convergence on \(X^{\divis}\) is Hausdorff.  If a
decreasing net \((\phi_i)\) in \(\PSH_X(\theta)\) converges pointwise on
\(X^{\divis}\) to the restriction of some
\(\phi\in\PSH_X(\theta)\), then \(\phi_i\searrow\phi\) pointwise on all of
\(X^{\An}\).
\end{theorem}

\begin{proof}
Finiteness and determination are
\Cref{thm:bj-basic-closure}\,(7), and Hausdorffness follows because the
divisorial evaluation maps separate psh functions.  The assertion about a
decreasing net is \cite[Corollary~4.25]{BJ22}: its pointwise limit on
\(X^{\An}\) is generically finite, since its values at the trivial
valuations are the corresponding finite values of \(\phi\).  It is therefore
\(\theta\)-psh by
\Cref{thm:bj-basic-closure}\,(4), has the prescribed divisorial values, and
therefore equals \(\phi\) by determination.
\end{proof}

\begin{warning}{Compactness is not unconditional}
\label{warn:ch34-compactness-conditional}
It is tempting to add to \Cref{thm:ch34-psh-compactness} the assertion that
the normalized space \(\PSH_X^{\sup}(\theta)\) of
\eqref{eq:normalized-psh-space} is \emph{compact} for this topology whenever
\(\theta\) is nef.  That assertion is false in general, and the error is not
subtle: by \Cref{thm:bj-envelope-equivalences} compactness of the normalized
space is \emph{equivalent} to the envelope property
\(\operatorname{EP}_X(\theta)\), and by \Cref{thm:ep-implies-unibranch} the
envelope property in an ample class forces \(X\) to be unibranch.  An
unconditional compactness theorem would therefore prove that every integral
projective variety is unibranch, which the nodal curve refutes.

The correct statement is conditional:
\emph{if \(\operatorname{EP}_X(\theta)\) holds, then
\(\PSH_X^{\sup}(\theta)\) is compact for the topology of pointwise
convergence on \(X^{\divis}\)}; this is
\Cref{thm:bj-envelope-equivalences}, and for \(X\) integral projective
unibranch the hypothesis is supplied by the envelope theorem,
\Cref{thm:bj-final-envelope}.  Because Chapters
\cref{chap:bj-cartier-realization}--\cref{chap:bj-subintegral-normalization}
prove that theorem, they may not use compactness in this form; each states
the prohibition explicitly, and the present chapter records it here so that
the unconditional version is never available to be misused.
\end{warning}

\section{Proper support and the value \texorpdfstring{\(-\infty\)}{-infinity}}
\label{sec:ch34-proper-support}

The class \(\PSH_X(\theta)\) contains functions taking the value
\(-\infty\), and the locus where this happens is governed entirely by
supports.  We record the mechanism, because it is the source of every
extended-real subtlety in Part~IV.

Let \(0\ne s\in H^0(X,mL)\).  By \eqref{eq:ch34-log-ideal} applied to the
ideal generated by \(s\),
\begin{equation}
\label{eq:divisor-section-pole}
 \log|s|(v)=-v(s),
 \qquad
 \log|s|(v)=-\infty
 \iff
 s(v)\subset\operatorname{div}(s),
\end{equation}
where \(s(v)\) is the support of \(v\).  In particular \(\log|s|\) is finite
at every point of \(X^{\val}\), since such a point has support an
irreducible component, on which \(s\) does not vanish identically; the
value \(-\infty\) occurs only at points of proper support.

\begin{example}
\label{ex:gf-proper-support-pole}
Let \(X=\mathbf P^1\), \(L=\cO(1)\), and let \(s\in H^0(X,L)\) be the
section vanishing exactly at the point \(p=0\).  The one-term system
\((1,(s),(0))\) presents
\(\phi:=\log|s|\in\cH^{\mathrm{gf}}(L)\); it is generically finite because
\(s\) does not vanish identically on \(X\), but it is not a Fubini--Study
function, since \(s\) has a zero.  Evaluating:
\(\phi(t\ord_0)=-t\) for \(t\ge0\); \(\phi(t\ord_q)=0\) for every closed
point \(q\ne0\) and \(t\ge0\); \(\phi(v_{X,\triv})=0\); and, at the trivial
semivaluation of the point \(0\),
\[
 \phi(v_{0,\triv})=-v_{0,\triv}(s)=-\infty ,
\]
because \(s\) vanishes at \(0\) and \(v_{0,\triv}\) assigns the value
\(+\infty\) to every function vanishing there.  So the pole locus of
\(\phi\) is the single point \(v_{0,\triv}\), which has proper support, in
accordance with \eqref{eq:divisor-section-pole}.
\end{example}

\begin{example}
\label{ex:ch34-decreasing-to-pole}
Keep the data of \Cref{ex:gf-proper-support-pole} and set
\[
 \phi_j:=\max\{\log|s|,\,-j\},\qquad j\in\mathbf Z_{>0}.
\]
Choose a base-point-free basis \(t_0,t_1\) of
\(H^0(\mathbf P^1,\cO(1))\).  Since
\(\max\{\log|t_0|,\log|t_1|\}=0\) on \(X^{\An}\), each \(\phi_j\) is the
genuine Fubini--Study function presented by
\[
                 \bigl(1,(s,t_0,t_1),(0,-j,-j)\bigr).
\]
The sequence is decreasing, and
\[
 \phi_j(v_{0,\triv})=\max\{-\infty,-j\}=-j\longrightarrow-\infty,
 \qquad
 \phi_j\downarrow\log|s| .
\]
Thus a decreasing sequence of \emph{finite-valued} Fubini--Study functions
converges to a function with a pole, and the pole appears only in the limit.
This is the elementary mechanism by which \(-\infty\) enters
\(\PSH_X(\theta)\), and it shows that no bound on the approximants can
prevent it.
\end{example}

We close with the smallest example in which the numerical constraint on a
psh function is visible, and which will be reused whenever a normalization
has to be checked.

Let \(X=\mathbf P^1\) and \(\theta=c_1(\cO(1))\), so \(\deg\theta=1\).  Let
\(s_0,s_1\in H^0(X,\cO(1))\) be the coordinate sections, vanishing at \(0\)
and at \(\infty\) respectively, and for \(c>0\) put
\begin{equation}
\label{eq:p1-fs}
 \phi_c:=\max\{\log|s_0|,\ \log|s_1|-c\}\in\cH_{\mathbf R}(\cO(1)).
\end{equation}

\begin{example}
\label{ex:p1-fs-function}
The pair \((s_0,s_1)\) has no common zero, so \(\phi_c\) is a Fubini--Study
function and is finite on all of \(X^{\An}\).  Its values are
\[
 \begin{aligned}
  \phi_c(t\ord_0)&=\max\{-t,\,-c\},\\
  \phi_c(t\ord_\infty)&=\max\{0,\,-t-c\}=0,\\
  \phi_c(t\ord_q)&=0
 \end{aligned}
\]
for every closed point \(q\notin\{0,\infty\}\), and
\(\phi_c(v_{X,\triv})=0\).  So \(\phi_c\) is the constant \(0\) except along
the ray through \(\ord_0\), where it decreases with slope \(-1\) until
\(t=c\) and is constant \(-c\) thereafter.  In particular \(\phi_c\) is
convex along every ray, with
\(\phi_c'{}_{,0}(0^+)=-1\) and \(\phi_c'{}_{,q}(0^+)=0\) for \(q\ne0\).

By the curve criterion recorded at the end of
\Cref{sec:ch34-pl-positivity}, \(\phi_c\) is \(\theta\)-psh if and only if
\(\deg\theta+\sum_q\phi_c'{}_{,q}(0^+)\ge0\), that is
\[
 1+(-1)=0\ \ge\ 0 ,
\]
which holds, with equality.  The example is therefore extremal: replacing
the coefficient \(1\) of \(\log|s_0|\) by any \(\lambda>1\) produces
\(1-\lambda<0\) and destroys pshness.  This is the smallest instance of the
principle that on a curve the total outgoing slope of a \(\theta\)-psh
function is bounded by the degree of the class, and it is the normalization
against which the energy formulae of \Cref{chap:bj-energy} are checked.
\end{example}

\section{Source note}

The definition of \(\PSH_X(\theta)\) and its varying-class decreasing
closure are \cite[Definition~4.1 and Theorem~4.5]{BJ22}; the diagonal
mechanism is \cite[Lemma~4.6]{BJ22}.  Constants, finite maxima, sums,
strictly positive scalar multiplication, pullback, and restriction to
components are \cite[Theorem~4.7]{BJ22}.  Notice that the scalar in the
source is strictly positive; no zero multiple of an extended-real function
is used.  Finiteness on \(X^{\lin}\) is
\cite[Corollary~4.17]{BJ22}, and divisorial comparison and determination are
\cite[Theorem~4.22 and Corollary~4.23]{BJ22}.
The three regularization scopes in \Cref{thm:ch34-fs-regularization} are
exactly \cite[Theorem~4.15\,(i)--(iii)]{BJ22}.

The closure block above remains a primary-source import without an
exact-current Rethlas receipt.  This provenance status is a gate against
uncited extensions, not a claim that the primary-source statements have
failed.

The later addendum \cite{BJ22b} retains this definition and closure package:
its new results concern the envelope property and the energy of a regularized
envelope.  In particular, it continues to invoke
\cite[Theorems~4.5 and~4.22]{BJ22}; it neither supplies nor requires a
different notion of decreasing closure, pullback, or divisorial
determination.
